\documentclass[11pt,a4paper]{article}
\usepackage[T1]{fontenc}
\usepackage[utf8]{inputenc}
\usepackage{lmodern}
\usepackage{microtype}
\usepackage[margin=27mm,headheight=14pt]{geometry}
\usepackage{amsmath,amssymb,amsthm,mathtools,mathrsfs}
\usepackage{booktabs,longtable,array,enumitem}
\usepackage[dvipsnames]{xcolor}
\usepackage{fancyhdr}
\usepackage{hyperref}
\hypersetup{colorlinks=true,linkcolor=MidnightBlue,citecolor=MidnightBlue,urlcolor=MidnightBlue,
 pdftitle={Surcomplex Analysis: Germs, Holomorphic Profiles, and Infinitesimal Zero Geometry},
 pdfauthor={OpenAI},pdfsubject={A constructive theory of analysis over No(i)}}
\usepackage[nameinlink,noabbrev]{cleveref}
\numberwithin{equation}{section}
\newtheorem{theorem}{Theorem}[section]
\newtheorem{proposition}[theorem]{Proposition}
\newtheorem{lemma}[theorem]{Lemma}
\newtheorem{corollary}[theorem]{Corollary}
\theoremstyle{definition}
\newtheorem{definition}[theorem]{Definition}
\newtheorem{example}[theorem]{Example}
\newtheorem{remark}[theorem]{Remark}
\newcommand{\No}{\mathbf{No}}
\newcommand{\K}{\mathbb K}
\newcommand{\C}{\mathbb C}
\newcommand{\R}{\mathbb R}
\newcommand{\N}{\mathbb N}
\newcommand{\Z}{\mathbb Z}
\newcommand{\tmon}{\mathfrak t}
\newcommand{\OO}{\mathcal O}
\newcommand{\mm}{\mathfrak m}
\newcommand{\HH}{\mathscr H}
\newcommand{\GG}{\mathscr G}
\newcommand{\PP}{\mathscr P}
\newcommand{\st}{\operatorname{st}}
\newcommand{\supp}{\operatorname{supp}}
\newcommand{\lc}{\operatorname{lc}}
\newcommand{\Res}{\operatorname{Res}}
\newcommand{\ord}{\operatorname{ord}}
\newcommand{\id}{\operatorname{id}}
\newcommand{\hhat}[1]{\widehat{#1}}
\newcommand{\hint}{\int^{\!H}}
\newcommand{\hoint}{\oint^{\!H}}
\newcommand{\vprof}{v_{\!H}}
\newcommand{\fin}{\mathrm{fin}}
\newcommand{\infp}{\uparrow}
\newcommand{\D}{\mathbb D}
\DeclareMathOperator{\Ree}{Re}
\DeclareMathOperator{\Imm}{Im}
\DeclareMathOperator{\divisor}{div}
\setlist{itemsep=3pt,topsep=5pt}
\setlength{\parskip}{3.3pt plus 1pt}
\setlength{\emergencystretch}{2em}
\pagestyle{fancy}
\fancyhf{}
\fancyhead[L]{\small Surcomplex analysis}
\fancyhead[R]{\small\nouppercase{\leftmark}}
\fancyfoot[C]{\thepage}
\renewcommand{\sectionmark}[1]{\markboth{Section \thesection}{}}
\renewcommand{\headrulewidth}{0.3pt}
\newcommand{\shortsection}[2]{\section[#1]{#2}\markboth{\thesection\enspace #1}{}}

\title{\textbf{Surcomplex Analysis}\\[6pt]
\large Hahn-analytic germs, holomorphic profiles,\\
and infinitesimal zero geometry}
\author{}
\date{21 September 2026}

\begin{document}
\maketitle
\begin{abstract}
We develop a two-level analytic theory over the algebraically closed proper-class
field $\K=\No(i)$. At the local level, every formal power series with surcomplex
coefficients becomes an actual function on a sufficiently small neighborhood,
when its sum is interpreted by Hahn summability rather than by convergence of
partial sums. We prove a scale-separation theorem, differentiation and
substitution rules, inverse and implicit function theorems, a local ramification
normal form, and local open-mapping and maximum-modulus principles.

At the global level, we introduce holomorphic profiles: Hahn series whose
coefficients are ordinary holomorphic functions on one common complex domain.
They lift canonically to its infinitesimal thickening. We prove an identity
theorem, Cauchy formulas at nonstandard points, a residue theorem, a constructive
Weierstrass preparation theorem, and an argument principle. In particular, a zero
of multiplicity $d$ of the leading complex coefficient lifts to exactly $d$
surcomplex zeros, counted with multiplicity, in its infinitesimal neighborhood.
Their power sums are recovered by contour integrals.

The distinction between the two levels is essential. In the full surreal
valuation topology every set-indexed convergent net is eventually constant;
Hahn sums need not be topological limits. We exhibit bounded nonconstant
analytic functions, nonunique globally analytic exponentials, and a global
locally analytic logarithm on $\K^\times$. These examples identify precisely why
classical global conclusions cannot be imported from local differentiability
alone. All assertions are made with explicit support, domain, and integration
conventions.
\end{abstract}
\thispagestyle{empty}
\clearpage
\tableofcontents
\clearpage

\section{The analytic structure has to be specified}

A surcomplex number is an expression $x+iy$ with $x,y\in\No$ and $i^2=-1$.
The ordered field $\No$ is real closed, so its complexification is algebraically
closed. This is an excellent algebraic starting point, but it does not by itself
supply a topology, a notion of summation, or a contour integral. Those choices
are part of an analytic theory, not consequences of algebraic closedness.
The foundational facts about $\No$ used here are standard; see
\cite{Gonshor,Alling,BMder}.

The aim is to retain both the infinitesimal richness of $\No$ and the
coherence of ordinary complex analysis. We use two complementary categories.

\paragraph{Local Hahn-analytic functions.}
Near an arbitrary $a\in\K$, including an infinite point, a function is represented
by a Hahn-summable evaluation of a power series
$\sum_{n\geq0}a_n(z-a)^n$. The coefficients may be arbitrary surcomplex numbers.
The resulting germ algebra is exactly $\K[[T]]$. Local algebra is therefore
unusually strong, while global uniqueness is unusually weak.

\paragraph{Holomorphic profiles.}
On an ordinary complex domain $U$, we consider a single controlled Hahn family
$F=\sum_\gamma f_\gamma\tmon^\gamma$, with every $f_\gamma$ holomorphic on $U$.
Its actual surcomplex arguments range over the entire infinitesimal thickening
$U^\sharp=\st^{-1}(U)$. The ordinary complex variable controls analytic
continuation and contours; the Hahn variable controls infinitesimal scales.
This additional coherence recovers a substantial global theory.

\subsection{The principal structural result}

The central zero-counting statement can be stated before the technical details.
Suppose
\[
 F(z)=f_0(z)+\sum_{\gamma\in S} f_\gamma(z)\tmon^\gamma,
 \qquad S\subset\No_{>0}\text{ a well-ordered set},
\]
and all coefficients are holomorphic on a common neighborhood of a complex
point $a$. If $f_0$ has a zero of multiplicity $d$ at $a$, then the lifted function
has exactly $d$ zeros, counted with multiplicity, with standard part $a$.
The proof constructs
\[
 F(a+T)=P(T)U(T),\qquad
 P(T)=T^d+p_{d-1}T^{d-1}+\cdots+p_0,\quad p_j\in\mm,
\]
where $U$ has nonvanishing leading complex coefficient. The construction uses
well-ordered recursion with finite coefficient calculations, not a completeness
assumption or a convergence assertion about Newton iterates.

This theorem links three kinds of information: an ordinary complex divisor,
a polynomial over the full surcomplex field, and infinitesimal positions of
actual zeros. The Cauchy and argument-principle formulas then recover the
positions through their symmetric functions.

\subsection{Relation to earlier work}

Surreal analysis is not a new subject. Alling's book develops formal power
series and hyperconvergence over surreal number fields \cite{Alling}; its
power-series chapter is directly relevant to the principle that arbitrary
formal coefficients can be accommodated by sufficiently small scales.
Rubinstein-Salzedo and Swaminathan study functions, limits, derivatives, and
integration with other explicitly chosen conventions \cite{RSS}.
Berarducci and Mantova construct surreal derivations and develop composition
and analyticity for substantial transseries classes \cite{BMder,BMgerms}.
The present variable derivative is not their derivation of the coefficient
field: every constant surcomplex coefficient has variable derivative zero.

There is also a developed theory called Hahn-holomorphic or Hahn-meromorphic
function theory, notably in work of M\"uller and Strohmaier \cite{MS}. Their
monomials are functions of a complex argument near a logarithmic covering
point and satisfy analytic convergence conditions. Here $\tmon^\gamma$ is a
surreal monomial independent of the ordinary variable, and the support may be
an arbitrary well-ordered \emph{set}. To avoid conflating these constructions,
we use the name \emph{holomorphic profile}.

The results below are a constructive synthesis with proofs in the stated
categories. The local formal results belong to the general tradition of
Hahn-series analysis, and the deformation arguments have the character of
Hensel and Weierstrass methods. No claim that these broad ideas originate here
is needed. The exact formulations, their compatibility, and the counterexamples
are what matter. Classical complex-analysis results used as inputs are
identified when invoked; accessible references include \cite{Orloff,ConwayComplex}.

\section{Algebra, valuation, and admissible sums}

\subsection{Class-size conventions and normal form}

We work in a class theory supporting the usual surreal construction, with
sets distinguished from proper classes. A coefficient sequence, a support,
a contour, and the family of data in any single construction are sets.
No sum in this article is indexed by all surreal numbers with proper-class
support. Recursion on a well-ordered support is set recursion. Topological
statements are understood through neighborhood predicates on classes; they do
not postulate a set of all open subclasses of $\K$.

Write a surcomplex normal form as
\begin{equation}\label{eq:normal}
 z=\sum_{\gamma\in S}c_\gamma\tmon^\gamma,
 \qquad c_\gamma\in\C,\qquad
 \tmon^\gamma:=\omega^{-\gamma},
\end{equation}
where $S$ is a well-ordered set in the additive ordered group $\No$ and the
power of $\omega$ is Conway's monomial operation. Thus
$\tmon^{\alpha+\beta}=\tmon^\alpha\tmon^\beta$.
This notation is \emph{not} a definition by the surreal exponential:
one must not replace $\log(\tmon^\gamma)$ by
$-\gamma\log\omega$ for arbitrary surreal $\gamma$.

For $z\neq0$, define
\[
 v(z)=\min\supp(z),\qquad \lc(z)=c_{v(z)},\qquad v(0)=+\infty.
\]
Then
\begin{equation}\label{eq:valrules}
 v(zw)=v(z)+v(w),\qquad
 v(z+w)\geq\min\{v(z),v(w)\}.
\end{equation}
Every nonzero element has a unique decomposition
\begin{equation}\label{eq:leadingunit}
 z=c\tmon^\gamma(1+\eta),\qquad
 c\in\C^\times,\quad\gamma\in\No,\quad v(\eta)>0.
\end{equation}
These are the usual Hahn-series descriptions of the surreal normal form
\cite{Alling,BMder}.

\subsection{Modulus and the residue field}

For $z=x+iy$, put
\[
 \bar z=x-iy,\qquad |z|=(x^2+y^2)^{1/2}\in\No_{\geq0}.
\]
Real closedness supplies the nonnegative square root. Algebraic
Cauchy--Schwarz gives
\[
 |zw|=|z||w|,\qquad |z+w|\leq|z|+|w|.
\]
The leading term of $|c\tmon^\gamma(1+\eta)|$ is
$|c|\tmon^\gamma$, so $v(|z|)=v(z)$.

Set
\[
 \OO=\{z:v(z)\geq0\},\qquad
 \mm=\{z:v(z)>0\}.
\]
These are the finite elements and the infinitesimals, respectively, with zero
included by the convention for $v(0)$. The standard-part map
\[
 \st:\OO\longrightarrow\C
\]
is the coefficient of $\tmon^0$. It is a surjective ring homomorphism with
kernel $\mm$. Hence $\OO/\mm\simeq\C$, and every finite $z$ is uniquely
$a+\eta$ with $a\in\C$ and $\eta\in\mm$.
For a set $U\subseteq\C$, write
\[
 U^\sharp=\{a+\eta:a\in U,\ \eta\in\mm\}.
\]
Even when $U$ is a set, $U^\sharp$ is generally a proper class.

\subsection{Strong or Hahn summability}

\begin{definition}\label{def:summable}
A set-indexed family $(z_j)_{j\in J}$ in $\K$ is \emph{Hahn summable}, or
\emph{strongly summable}, if
\[
 \bigcup_{j\in J}\supp(z_j)
\]
is well ordered and each exponent belongs to the supports of only finitely
many members of the family. Its sum is formed by adding, at each exponent,
that finite collection of complex coefficients.
\end{definition}

Thus all coefficients of a Hahn sum are obtained by finite arithmetic.
A summable family can be regrouped arbitrarily, and a subfamily is summable.
Products of summable families can be expanded and regrouped whenever the
combined family meets the same support criterion. This is a summation
operation, not an assertion about limits of finite partial sums.

\begin{lemma}[Neumann's summability lemma]\label{lem:neumann}
Let $A,B$ be well-ordered subsets of an ordered abelian group. Then $A+B$ is
well ordered, and each element of the group has only finitely many
representations $a+b$ with $a\in A$, $b\in B$. If $S$ is a well-ordered subset
of the strictly positive cone, its generated additive monoid
\[
 M(S)=\{s_1+\cdots+s_n:n\geq0,\ s_j\in S\}
\]
is well ordered, and each group element has only finitely many representations
as an ordered finite word in $S$.
\end{lemma}

The positivity in the last assertion is essential. This standard
combinatorial lemma underlies Hahn fields and geometric-series inversion;
see \cite{Neumann,MS}. For the two-set assertion, a decreasing sequence of sums
would have a subsequence with nondecreasing first coordinates, forcing its
second coordinates to decrease. The same argument rules out infinitely many
decompositions of a fixed sum. The finite-word assertion is the corresponding
Neumann lemma for arbitrary word lengths; it does not require an Archimedean
value group.

\begin{corollary}\label{cor:infseries}
If $\eta_1,\ldots,\eta_r\in\mm$ and
$A\in\C[[X_1,\ldots,X_r]]$, then $A(\eta_1,\ldots,\eta_r)$ is a well-defined
Hahn sum. Formal identities, differentiation identities, and compositions
with zero constant term may be evaluated this way.
\end{corollary}
\begin{proof}
The union of the finitely many positive supports is well ordered. Every
expanded monomial has support in its generated monoid. Neumann's lemma gives
both well-ordering and the finite number of words that can contribute to
any exponent. There are only finitely many assignments of a word to the
fixed number of variables. Hence all formal expansions in question are
legitimate regroupings of summable families.
\end{proof}

In particular,
\begin{equation}\label{eq:small-explog}
 \exp_H(\eta)=\sum_{n\geq0}\frac{\eta^n}{n!},\qquad
 \log_H(1+\eta)=\sum_{n\geq1}\frac{(-1)^{n+1}\eta^n}{n}
\end{equation}
are inverse group isomorphisms between $(\mm,+)$ and $(1+\mm,\cdot)$.
The same convention defines $(1+\eta)^\alpha$ for $\alpha\in\C$ by its formal
binomial series. In particular, these operations commute with conjugation.

\section{What the full surreal topology does and does not do}

We use the topology with basic neighborhoods
\[
 B_\rho(a)=\{z:v(z-a)>\rho\},\qquad\rho\in\No.
\]
We call it the \emph{sharp topology}. It is the same as the topology generated
by the surreal-modulus balls $|z-a|<r$, $r\in\No_{>0}$. To check the
coincidence, a sufficiently large valuation makes a quantity smaller than a
given $r$, while a ball of radius $\tmon^{\rho+1}$ is contained in $B_\rho(0)$.
Every valuation ball is also closed, by the ultrametric inequality.

The following elementary size fact controls the topology.

\begin{lemma}[A set has room above and below it]\label{lem:room}
Every set of surreal numbers is bounded above and below by surreal numbers.
Every set $E\subset\No_{>0}$ has a strictly positive lower bound strictly
smaller than every member of $E$.
\end{lemma}
\begin{proof}
For a set $A$, the surreal cuts $\{A\mid\}$ and $\{\mid A\}$ provide an
upper and a lower bound. For the second assertion the cut $\{0\mid E\}$ is a
positive surreal smaller than every member of $E$.
\end{proof}

\begin{theorem}[Set-discreteness]\label{thm:setdiscrete}
Every set $A\subset\K$ is closed and discrete in the induced sharp topology.
Every convergent net in $\K$ indexed by a set is eventually equal to its limit.
In particular, every sharply convergent ordinary sequence is eventually constant.
\end{theorem}
\begin{proof}
For any $p\in\K$, apply \cref{lem:room} to the set of nonzero distances
$\{|p-a|:a\in A,\ a\neq p\}$. A sufficiently small modulus ball about $p$
contains no point of $A$ other than possibly $p$ itself. This proves both
closedness and discreteness.

If $(z_j)_{j\in J}$ converges to $z$ and $J$ is a set, the nonzero distances
$|z_j-z|$ also form a set. A ball smaller than all these distances contains no
value of the net except $z$. Convergence therefore makes the net eventually
identically $z$.
\end{proof}

\begin{example}[A sum which is not a sequential limit]\label{ex:geom}
For $t=\tmon^1=\omega^{-1}$, Neumann's lemma gives
\[
 \sum_{n\geq0}t^n=(1-t)^{-1}.
\]
But the remainder after degree $N$ has valuation $N+1$, which never exceeds
$\omega$. Thus the partial sums do not converge sharply to $(1-t)^{-1}$.
Even the real sequence $1/n$ does not converge sharply to zero.
\end{example}

This is not a defect in the sum. It is a warning against giving one operation
the name of another. Strong summation is coefficientwise and algebraic;
sharp convergence tests all surreal precision scales at once.

\begin{proposition}[Ordinary paths are not sharp paths]\label{prop:paths}
A continuous map from a connected ordinary topological space into $\K$ whose
image is a set is constant. In particular, a nonconstant usual complex contour,
viewed as a map from a real interval into $\K$, is not sharp-continuous.
\end{proposition}
\begin{proof}
Its image has the discrete subspace topology by \cref{thm:setdiscrete}.
A continuous image of a connected space is connected, and a discrete space
has no connected subset containing two points.
\end{proof}

Consequently, our later contour integral will be explicitly \emph{coefficientwise
along an ordinary complex contour}. It is not an unspecified surreal-valued
Riemann integral. Other topologies and transfinite notions of convergence may
be useful, but their theorems cannot be silently mixed with this topology.

\begin{example}[Local constancy is not global constancy]\label{ex:indicator}
The function on all of $\K$ that is zero on $\OO$ and one on $\K\setminus\OO$
is locally constant, hence locally power-series analytic, with derivative zero.
It is bounded and nonconstant. The set $\OO$ and its complement are both open.
This already rules out an unrestricted global identity theorem, Liouville
theorem, or uniqueness theorem for a differential equation based only on
local power-series analyticity.
\end{example}

\section{Every formal series defines a local analytic germ}

\subsection{A scale-separation theorem}

The complex version of the small-radius phenomenon is particularly transparent
with the full surreal value group. The essential point is to bound
\emph{supports}, not merely leading valuations.

\begin{theorem}[Scale separation]\label{thm:scale}
Let $(a_n)_{n\geq0}$ be any sequence in $\K$. There exists $\beta>0$ such that
\[
 \supp(a_n)\subset[-\beta,\beta]\quad\text{for every }n.
\]
For every $h\in\K$ with $v(h)>2\beta$, the family $(a_nh^n)_{n\geq0}$ is
Hahn summable. Its sum defines a function on $B_{2\beta}(0)$.
The analogous assertion holds for formal series in finitely many variables,
with $\min_jv(h_j)>2\beta$.
\end{theorem}
\begin{proof}
The union of the coefficient supports is a set, even though it need not be
well ordered. By \cref{lem:room} it lies in some $[-\beta,\beta]$.
The case $h=0$ is immediate. Otherwise write
\[
 h=c\tmon^\lambda(1+u),\qquad\lambda>2\beta,\quad u\in\mm.
\]
The sets
\[
 A_n=\supp(a_n)+n\lambda\subset[n\lambda-\beta,n\lambda+\beta]
\]
are individually well ordered and occur in strictly separated increasing
bands. Their union $A$ is therefore well ordered, and its elements determine
$n$ uniquely. Put $M=M(\supp u)$. Expanding the finite polynomial $(1+u)^n$
shows that all the supports of $a_nh^n$ lie in $A+M$.

For a fixed exponent, Neumann's lemma gives finitely many decompositions into
an element of $A$ and an element of $M$. The former fixes $n$; the latter has
finitely many contributing words in $\supp u$. This proves local finiteness
as well as well-ordering, hence summability.

For finitely many nonzero $h_j$, put $\lambda=\min_jv(h_j)$ and write
$h_j=\tmon^\lambda u_j$ with $u_j\in\OO$. The union of the strictly positive
supports of these finitely many $u_j$ is well ordered. Group the formal
monomials by total degree $n$, of which there are finitely many in each degree.
The same separated bands and generated positive monoid prove the assertion.
\end{proof}

\begin{remark}
The proof does not require $n\lambda$ to tend to infinity in the full value
group. It usually does not. It only requires separation of the coefficient
bands and the finite-word property.
\end{remark}

\begin{example}[Valuation bounds alone are insufficient]\label{ex:badcoeff}
For $n\geq1$, take
\[
 a_n=\tmon^{\omega-n+1/n}.
\]
Every $a_n$ is infinitesimal. Nevertheless, at $h=\tmon$ the proposed summands
are $a_nh^n=\tmon^{\omega+1/n}$, whose support is strictly decreasing and
has no least element. The family is not summable. An assertion that
``all finite coefficients can be evaluated at every infinitesimal'' would be
false. The common-support hypotheses used later prevent precisely this problem.
\end{example}

\begin{example}[A radius unavailable at rank one]
The formal series $\sum_{n\geq0}\tmon^{-n^2}T^n$ cannot be strongly evaluated
at a nonzero argument of positive real valuation: the valuations
$-n^2+n\lambda$ are unbounded below. Over the full surcomplex field it is
summable at $T=\tmon^\omega$, since $n\omega-n^2$ is strictly increasing.
More generally, \cref{thm:scale} supplies a whole neighborhood.
\end{example}

\subsection{Differentiation and recentering}

A derivative will always be a full sharp limit, equivalently a
surreal-modulus $\varepsilon$--$\delta$ limit. The domain of a difference quotient
contains a proper class of small increments, so \cref{thm:setdiscrete} does not
make this definition vacuous.

\begin{theorem}[Taylor calculus]\label{thm:taylor}
Let $A(T)=\sum_{n\geq0}a_nT^n$, and choose $\beta$ as in
\cref{thm:scale}. On $B_{2\beta}(0)$ its Hahn evaluation is differentiable and
\[
 A'(h)=\sum_{n\geq1}na_nh^{n-1}.
\]
It has derivatives of all finite orders. Whenever $h,k\in B_{2\beta}(0)$,
its binomial expansion may be regrouped as
\begin{equation}\label{eq:taylor-shift}
 A(h+k)=\sum_{j\geq0}\frac{A^{(j)}(h)}{j!}k^j.
\end{equation}
The displayed family is Hahn summable. In particular, for sufficiently small
nonzero $k$,
\begin{equation}\label{eq:remainder}
 v\bigl(A(h+k)-A(h)-A'(h)k\bigr)\geq2v(k)-\beta.
\end{equation}
\end{theorem}
\begin{proof}
Expand $a_n(h+k)^n$ as a sum over $0\leq j\leq n$. This is a series in two
variables with finitely many terms in each total degree; the proof of
\cref{thm:scale} applies to the fully expanded family. Regrouping first in $j$
gives \eqref{eq:taylor-shift}, with the indicated derivative coefficients.

In the tail with $j\geq2$, each term has valuation at least
\[
 -\beta+(n-j)v(h)+jv(k)\geq-\beta+2v(k).
\]
A Hahn sum of terms with a common lower valuation bound has the same lower
bound, proving \eqref{eq:remainder}. After division by $k$, the error has
valuation at least $v(k)-\beta$. Given any target valuation $\rho$, choose
$v(k)>\rho+\beta$ as well as $v(k)>2\beta$. This proves the derivative claim.
The same argument applies repeatedly; multiplication of coefficients by
ordinary integers does not enlarge their supports.
\end{proof}

\begin{lemma}[Formal identities pass to germs]\label{lem:formal-eval}
Addition, multiplication, formal differentiation, and composition
$A(B(T))$ with $B(0)=0$ commute with Hahn evaluation after shrinking the
neighborhood if necessary. The assertion also holds in finitely many variables.
\end{lemma}
\begin{proof}
For each total degree, every coefficient in a formal product or such a
composition is a finite sum of finite products of input coefficients.
Consider all the individual coefficient products before collecting like
terms, over all degrees. This is a set of surcomplex numbers. Their supports
have a common surreal bound by \cref{lem:room}. The separated-band proof
therefore applies to the fully expanded family, with finitely many
combinatorial terms in each degree. All parenthesizations and groupings give
the same sum. Shrinking also ensures that the evaluated inner series lies
in a neighborhood on which the outer series is defined, since its constant
term is zero and its evaluation is continuous.
\end{proof}

\subsection{The germ algebra and its local theorems}

\begin{definition}
A function on a sharp-open class is \emph{germ-analytic} if near each point
$a$ it agrees with a Hahn evaluation of a formal power series in $z-a$.
Denote its germ algebra at $a$ by $\GG_a$.
\end{definition}

This definition deliberately requires a power-series representation.
We do not assert a surcomplex Goursat theorem saying that the existence of
one complex derivative everywhere implies this representation.

\begin{theorem}[Classification of analytic germs]\label{thm:germs}
Hahn evaluation gives an isomorphism of $\K$-algebras
\[
 \K[[T]]\xrightarrow{\ \sim\ }\GG_a,
 \qquad \sum a_nT^n\longmapsto\left[z\mapsto\sum a_n(z-a)^n\right].
\]
The coefficient of $T^n$ is the $n$th derivative at $a$ divided by $n!$.
Every analytic germ has a primitive. A germ is a unit exactly when its
constant coefficient is nonzero.
\end{theorem}
\begin{proof}
Existence and compatibility are \cref{thm:scale,lem:formal-eval}.
Surjectivity is the definition of $\GG_a$. For injectivity, let $m$ be the
first index with $a_m\neq0$. If $v(h)$ is sufficiently large, the tail after
$a_mh^m$ has valuation at least $(m+1)v(h)-\beta$, strictly greater than
$v(a_m)+mv(h)$. Thus the value is nonzero for all sufficiently small nonzero
$h$. A nonzero formal series cannot represent the zero germ.

The derivative formula gives the coefficients. Formal integration divides
$a_n$ by $n+1$ and shifts degree; its evaluation is a primitive.
Formal inversion applies exactly to nonzero constant terms, and
\cref{lem:formal-eval} evaluates the inverse identity.
\end{proof}

\begin{theorem}[Inverse and implicit function theorems]\label{thm:inverse}
If $f$ is germ-analytic at $a$ and $f'(a)\neq0$, then it has a germ-analytic
inverse between suitable sharp neighborhoods of $a$ and $f(a)$.
More generally, if a germ-analytic map $F(X,Y)$ in finitely many variables
satisfies $F(0,0)=0$ and its square Jacobian matrix with respect to $Y$ at
$(0,0)$ is invertible, there is a unique formal, and hence germ-analytic,
solution $Y=G(X)$ with $G(0)=0$.
\end{theorem}
\begin{proof}
For one variable translate to $a=f(a)=0$, and write
$f(T)=a_1T+a_2T^2+\cdots$ with $a_1\neq0$. In
$g(S)=b_1S+b_2S^2+\cdots$, the equation $f(g(S))=S$ determines
$b_1=a_1^{-1}$ and then $b_n$ uniquely, because the degree-$n$ equation has
linear term $a_1b_n$ and otherwise involves earlier $b_j$. Substitution shows
$g(f(T))=T$ as well, either by the same uniqueness or by formal cancellation.
The identities evaluate on sufficiently small neighborhoods. Restricting
one neighborhood to the inverse image of the other gives an actual local
bijection.

In several variables, let $J=\partial_YF(0,0)$. Inductively suppose all
homogeneous terms of $G$ below degree $n$ are known. The degree-$n$ equation
has the form $JG_n+R_n=0$, where $R_n$ is already known, so
$G_n=-J^{-1}R_n$. This proves existence and uniqueness by degree recursion.
The multivariable scale theorem evaluates the solution and its formal identity.
Equivalently apply the formal inverse construction to $(X,Y)\mapsto(X,F(X,Y))$;
this also gives local uniqueness among actual sufficiently small solutions.
\end{proof}

\begin{theorem}[Ramification normal form]\label{thm:normalform}
If $f-f(a)$ is a nonzero analytic germ of order $m\geq1$, there is an invertible
analytic coordinate $\psi$ at $a$, with $\psi(a)=0$, such that
\begin{equation}\label{eq:ramification}
 f(z)=f(a)+\psi(z)^m.
\end{equation}
For all sufficiently small nonzero $w-f(a)$ there are exactly $m$ distinct
solutions of $f(z)=w$ in an appropriate fixed neighborhood of $a$.
At $w=f(a)$ there is the single zero $a$ of multiplicity $m$ in that neighborhood.
\end{theorem}
\begin{proof}
Write $f(a+T)-f(a)=T^m u(T)$ with $u(0)\neq0$. Choose an $m$th root
$c$ of $u(0)$ in $\K$. The formal binomial series constructs
$b(T)=c\bigl(u(T)/u(0)\bigr)^{1/m}$, with $b(T)^m=u(T)$.
The coordinate $Tb(T)$ has nonzero linear coefficient and is invertible by
\cref{thm:inverse}. Evaluation proves \eqref{eq:ramification}.

Choose a valuation ball in the image of this coordinate and restrict the
source to its inverse image. The polynomial $Y^m-s$ has exactly $m$ distinct
roots when $s\neq0$, all of valuation $v(s)/m$. For sufficiently large $v(s)$,
all lie in that image ball, and there are no others there. For $s=0$ the
only root is zero, with multiplicity $m$.
\end{proof}

\begin{corollary}[Open mapping and maximum modulus]\label{cor:openmax}
At a point where its germ is nonconstant, a germ-analytic function maps every
sufficiently small neighborhood onto a set containing a neighborhood of its
value. A germ-analytic function whose germ is nonconstant at every point is
an open map. Such a function has no local maximum of its surreal modulus.
Its real part has no local maximum or local minimum either.
\end{corollary}
\begin{proof}
The power map in \eqref{eq:ramification} maps a valuation ball onto a valuation
ball, using algebraic closedness and $v(Y^m)=mv(Y)$. The invertible coordinate
is a local homeomorphism. This proves openness. Every neighborhood of a
nonzero $b$ contains $(1+\delta)b$ for some positive infinitesimal $\delta$,
which has larger modulus; neighborhoods of zero contain nonzero points.
Similarly, any neighborhood of $b$ contains points with larger and smaller
real part. None of these extrema is therefore possible.
\end{proof}

The modulus in this corollary is $|z|=\sqrt{z\bar z}$, not an ultrametric
absolute value obtained from $v$. A valuation can be locally constant on
nonzero values, and would not satisfy the same maximum principle.

\subsection{Cauchy--Riemann equations and harmonic components}

Writing $f=u+iv$ and a small increment as $h=\xi+i\eta$, the first-order
term $f'(a)h$ gives
\[
 u_x=v_y,\qquad u_y=-v_x.
\]
The same power-series calculus justifies repeated real partial derivatives,
so
\[
 u_{xx}+u_{yy}=0,\qquad v_{xx}+v_{yy}=0.
\]
These are genuine local analogs of the Cauchy--Riemann and harmonicity
statements. Conversely, a real Fr\'echet derivative satisfying the
Cauchy--Riemann equations is $\K$-linear and hence gives a complex derivative.
This converse is a statement about the derivative, not an unproved
characterization of germ-analyticity.

\section{Holomorphic profiles and their canonical lifts}

\subsection{The profile algebra}

\begin{definition}\label{def:profile}
For an ordinary complex domain $U$, a \emph{holomorphic profile} is a formal
Hahn series
\begin{equation}\label{eq:profile}
 F=\sum_{\gamma\in S}f_\gamma\tmon^\gamma,
 \qquad f_\gamma\in\OO(U),
\end{equation}
where the set $S=\{\gamma:f_\gamma\not\equiv0\}$ is well ordered. Denote the
algebra of these profiles by $\HH(U)$. Its valuation is
$\vprof(F)=\min S$ for $F\neq0$, with $\vprof(0)=+\infty$. A profile of strictly positive valuation is
an \emph{infinitesimal profile}.
\end{definition}

Here $\OO(U)$ means ordinary complex holomorphic functions; it should not be
confused with the previously defined finite-element ring $\OO$. The context
and the presence of a domain distinguish them. A profile is not merely an
arbitrary function $U\to\K$: its coefficients must be holomorphic on one
common domain and its global support must be well ordered.

Neumann's lemma makes $\HH(U)$ a $\K$-algebra under coefficientwise addition
and Hahn multiplication. Constants from all of $\K$, including infinite ones,
are allowed. Define $F'=\sum f_\gamma'\tmon^\gamma$. Differentiation commutes
with algebra operations, and $\vprof(F')\geq\vprof(F)$.

This globally controlled algebra is naturally a presheaf under restriction.
We do not assume that arbitrary infinite gluing preserves a single global
well-ordered support. A sheaf version can instead require such support control
locally, but the global theorems below are stated in $\HH(U)$ itself.

\begin{theorem}[Canonical lift]\label{thm:lift}
For $F\in\HH(U)$ and $z=a+\eta\in U^\sharp$, the expression
\begin{equation}\label{eq:lift}
 \hhat F(a+\eta)=
 \sum_{\gamma\in S}\sum_{n\geq0}
 \frac{f_\gamma^{(n)}(a)}{n!}\tmon^\gamma\eta^n
\end{equation}
is Hahn summable. It defines a germ-analytic function on $U^\sharp$, and
\[
 \hhat{F+G}=\hhat F+\hhat G,\qquad
 \hhat{FG}=\hhat F\hhat G,\qquad
 (\hhat F)'=\hhat{F'}.
\]
For every $z\in U^\sharp$ and $h\in\mm$, one has the summable Taylor identity
\begin{equation}\label{eq:profile-taylor}
 \hhat F(z+h)=\sum_{n\geq0}\frac{\hhat{F^{(n)}}(z)}{n!}h^n.
\end{equation}
\end{theorem}
\begin{proof}
Let $M$ be the positive monoid generated by $\supp\eta$.
Every expanded term in \eqref{eq:lift} lies in $S+M$. For a fixed exponent,
there are finitely many decompositions into an element of $S$ and a word
from $\supp\eta$. The word length is $n$, so only finitely many terms of the
double family contribute. This proves summability.

The ordinary Taylor identities for addition and multiplication can now be
regrouped within these summable families. For \eqref{eq:profile-taylor},
expand at the common standard part $a=\st z$, using the two infinitesimals
$\eta=z-a$ and $h$. The positive monoid generated by their joint support again
provides finite coefficient calculations, and the binomial formula gives
the displayed identity. If $\gamma_0=\min S$, the tail after the linear term
has valuation at least $\gamma_0+2v(h)$. Division by $h$ proves the derivative
formula in the sharp topology. The Taylor identity also supplies a local
power-series representation, so the lift is germ-analytic.
\end{proof}

\begin{corollary}[Leading coefficient and standard part]\label{cor:reduction}
If $F\neq0$ has leading exponent $\gamma_0$ and leading coefficient
$f_{\gamma_0}$, then for $z=a+\eta\in U^\sharp$,
\[
 v\bigl(\hhat F(z)\bigr)\geq\gamma_0,
 \qquad
 \st\bigl(\tmon^{-\gamma_0}\hhat F(z)\bigr)=f_{\gamma_0}(a).
\]
In particular, a zero of $\hhat F$ can occur only in a monad above a zero
of $f_{\gamma_0}$.
\end{corollary}
\begin{proof}
After factoring $\tmon^{\gamma_0}$, all terms except the constant Taylor
term of $f_{\gamma_0}$ have strictly positive valuation.
\end{proof}

\subsection{Identity and uniqueness of continuation}

\begin{theorem}[Identity theorem for profiles]\label{thm:profileidentity}
Let $U$ be connected and $F,G\in\HH(U)$. Each of the following implies $F=G$:
\begin{enumerate}[label=(\roman*)]
\item their values agree on a set of ordinary complex points with an
accumulation point in $U$;
\item their lifts agree on a nonempty sharp-open subset of $U^\sharp$;
\item all derivatives of their lifts agree at one point of $U^\sharp$.
\end{enumerate}
In particular, a nonconstant profile on a connected domain has a nonconstant
analytic germ at every point of its thickening.
\end{theorem}
\begin{proof}
Replace $F-G$ by $F$. At an ordinary complex point $a$, normal-form uniqueness
says that $\hhat F(a)=0$ implies $f_\gamma(a)=0$ for every $\gamma$.
The classical identity theorem applied coefficientwise proves (i).

For (iii), let $z=a+\eta$ be the point in question. Applying
\eqref{eq:profile-taylor} to every derivative, with increment $-\eta$, gives
\[
 \hhat{F^{(j)}}(a)=
 \sum_{n\geq0}\frac{\hhat{F^{(j+n)}}(z)}{n!}(-\eta)^n=0.
\]
Normal-form uniqueness implies $f_\gamma^{(j)}(a)=0$ for all $j,\gamma$.
Ordinary holomorphic uniqueness then gives $f_\gamma\equiv0$ on $U$.
Agreement on a sharp-open subset gives equality of all derivatives at one
of its points, so (ii) follows from (iii).

Finally, if $\hhat F$ has a constant germ, its derivative has zero germ.
Applying the proven result to $F'$ shows $F'=0$, and connectedness makes all
its coefficient functions constant.
\end{proof}

Thus global uniqueness is recovered not from sharp connectedness, which is
unavailable, but from coherence of the ordinary coefficient functions.
Combining this theorem with \cref{cor:openmax} gives an open-mapping and
maximum-modulus principle for every nonconstant profile on a connected base.

\subsection{Substitution of infinitesimal deformations}

\begin{theorem}[Profile composition]\label{thm:profilecomposition}
Let $g:V\to U$ be ordinary holomorphic, $H\in\HH(V)$ have strictly positive
valuation, and $F\in\HH(U)$. Then
\begin{equation}\label{eq:profilecomp}
 F\circ(g+H):=
 \sum_{\substack{\gamma\in\supp F\\n\geq0}}
 \frac{f_\gamma^{(n)}(g)}{n!}\tmon^\gamma H^n
\end{equation}
defines a member of $\HH(V)$. Its lift is the actual composition
$\hhat F\circ\hhat{(g+H)}$. The ordinary chain rule holds.
\end{theorem}
\begin{proof}
If $S$ is the support of $F$ and $A$ is the positive support of $H$, the
expanded profile has support in $S+M(A)$. At each exponent only finitely many
words contribute. Their coefficient functions are finite sums of products
of holomorphic functions on $V$, so they are holomorphic there.
Evaluation at $b+\eta$ adds the positive support of $\eta$ to this same
finite-word calculation. The ordinary formal Taylor composition identity
then proves the lift assertion. Differentiating the formal expression, with
all regroupings justified by the same support set, proves the chain rule.
\end{proof}

For example, if $H$ has positive profile valuation, then
$\exp_H(H)$ and $\log_H(1+H)$ are profiles. If
$F=f_0+H$ with $f_0$ nowhere zero and $H$ of positive valuation, then
\begin{equation}\label{eq:profileinverse}
 F^{-1}=f_0^{-1}\sum_{n\geq0}(-H/f_0)^n\in\HH(U).
\end{equation}
This is an inverse in the profile algebra and, after lifting, a pointwise
inverse on all of $U^\sharp$.

\section{Cauchy theory with a specified contour integral}

\subsection{Coefficientwise integration}

Let $\Gamma$ be an ordinary piecewise $C^1$ complex path with compact image.
Suppose $F=\sum f_\gamma\tmon^\gamma$ has one well-ordered support and each
coefficient is continuous on the path. Define
\begin{equation}\label{eq:integral}
 \hint_\Gamma F(\zeta)\,d\zeta
 :=\sum_\gamma\tmon^\gamma
     \int_\Gamma f_\gamma(\zeta)\,d\zeta.
\end{equation}
The outer sum is a Hahn sum whose support is a subset of the original support;
the inner integrals are the ordinary complex integrals. The notation $H$
records this convention.

This integral is $\K$-linear and invariant under the usual piecewise smooth
reparametrizations. Both assertions follow coefficientwise; multiplication
by a surcomplex constant uses a Hahn product with finite coefficients.
For profiles, the usual fundamental-theorem identity at ordinary endpoints is
\[
 \hint_\Gamma F'(\zeta)\,d\zeta
   =\hhat F(\Gamma(1))-\hhat F(\Gamma(0)).
\]
No claim is made here about an arbitrary sharp path or an arbitrary
surreal-valued Riemann sum.

\begin{theorem}[Cauchy's theorem and primitives]\label{thm:cauchythm}
If $F\in\HH(U)$ and an ordinary closed contour $\Gamma$ is null-homologous
in $U$, then
\[
 \hoint_\Gamma F(\zeta)\,d\zeta=0.
\]
If $U$ is simply connected, every $F\in\HH(U)$ has a primitive in $\HH(U)$.
\end{theorem}
\begin{proof}
Classical Cauchy's theorem makes the integral of each coefficient zero.
For a primitive, choose an ordinary base point $a_0\in U$ and let
$g_\gamma(z)=\int_{a_0}^z f_\gamma(\zeta)\,d\zeta$, using the classical
path-independent integral. The family $\sum g_\gamma\tmon^\gamma$ has support
contained in the original one and derivative $F$.
\end{proof}

\begin{theorem}[Cauchy formula at a surcomplex point]\label{thm:cauchyformula}
Let $\Gamma$ be a positively oriented piecewise smooth Jordan contour, with
interior $D$, and suppose $F\in\HH(U)$ for a neighborhood $U$ of $\overline D$.
For $w=a+\eta\in D^\sharp$ and every integer $m\geq0$,
\begin{equation}\label{eq:cauchysharp}
 \hhat{F^{(m)}}(w)=
 \frac{m!}{2\pi i}\hoint_\Gamma
       \frac{F(\zeta)}{(\zeta-w)^{m+1}}\,d\zeta.
\end{equation}
The integrand is defined by expanding its kernel about $a$ and taking the
resulting Hahn profile along the contour.
\end{theorem}
\begin{proof}
For $\zeta$ near the contour, $\zeta-a\neq0$. The valid Hahn expansion is
\[
 \frac1{(\zeta-a-\eta)^{m+1}}
 =\sum_{n\geq0}\binom{m+n}{n}
       \frac{\eta^n}{(\zeta-a)^{m+n+1}}.
\]
Its product with $F$ is supported in $S+M(\supp\eta)$, with finite
contributions at every exponent. Therefore \eqref{eq:integral} applies.
Classical Cauchy differentiation gives
\[
 \frac{m!}{2\pi i}\int_\Gamma
 \binom{m+n}{n}\frac{f_\gamma(\zeta)}{(\zeta-a)^{m+n+1}}\,d\zeta
 =\frac{f_\gamma^{(m+n)}(a)}{n!}.
\]
The resulting double Hahn sum is exactly the lift of $F^{(m)}$ at $a+\eta$.
\end{proof}

This formula determines values at nonstandard points from a single ordinary
contour. It is stronger than merely writing down the classical formula at
ordinary complex arguments, but its validity depends on the controlled
profile structure.

\subsection{Morera and coefficientwise Cauchy estimates}

\begin{proposition}[Morera's theorem for profiles]
Let $F=\sum f_\gamma\tmon^\gamma$ have a common well-ordered support and
continuous complex coefficients on $U$. If its coefficientwise integral
around every triangle contained in $U$ is zero, then $F\in\HH(U)$.
\end{proposition}
\begin{proof}
Normal-form uniqueness says that the ordinary integral of each $f_\gamma$
around every such triangle is zero. Classical Morera's theorem makes every
coefficient holomorphic.
\end{proof}

If $\overline{D(a,r)}\subset U$, define the real number
$M_\gamma=\max_{|\zeta-a|=r}|f_\gamma(\zeta)|$. Then the ordinary Cauchy
estimate gives, for every $n,\gamma$,
\begin{equation}\label{eq:cauchyest}
 |f_\gamma^{(n)}(a)|\leq \frac{n!M_\gamma}{r^n}.
\end{equation}
These are simultaneous estimates at every Hahn scale. No uniform bound on
the numbers $M_\gamma$ is needed for the profile calculus: each appears as a
complex coefficient at its own exponent. On the other hand, a single
surreal pointwise bound is generally too coarse to recover these separate
bounds. This distinction will explain the failure of a naive Liouville theorem.

\subsection{Profile residues}

\begin{definition}
A \emph{meromorphic profile with common pole set $E$} is a Hahn series
$F=\sum f_\gamma\tmon^\gamma$ with common well-ordered support, where each
$f_\gamma$ is meromorphic on $U$ and all its poles lie in the same closed
discrete subset $E\subset U$. At $a\in E$, define its \emph{profile residue}
by
\[
 \Res^H_a F=\sum_\gamma
       \bigl(\Res_{z=a}f_\gamma(z)\bigr)\tmon^\gamma.
\]
There is no requirement that pole orders be uniformly bounded in $\gamma$.
\end{definition}

\begin{theorem}[Residue theorem]\label{thm:residue}
Let $\Gamma$ be an ordinary positively oriented Jordan contour in $U\setminus E$
whose closed interior lies in $U$. Then
\begin{equation}\label{eq:residue}
 \hoint_\Gamma F(\zeta)\,d\zeta
   =2\pi i\sum_{a\in E\cap\operatorname{int}\Gamma}\Res^H_a F.
\end{equation}
\end{theorem}
\begin{proof}
Only finitely many points of the closed discrete set $E$ lie in the compact
closed interior. Apply the classical residue theorem to every coefficient
and interchange a finite sum with the Hahn sum.
\end{proof}

The qualification ``profile'' matters. Consider a nonzero infinitesimal
$\varepsilon$. On an ordinary annulus about zero,
\begin{equation}\label{eq:movingpole}
 \frac1{z-\varepsilon}
   =\sum_{n\geq0}\varepsilon^n z^{-n-1}.
\end{equation}
After expanding the powers of $\varepsilon$, this is a meromorphic profile
with common pole set $\{0\}$, and its profile residue at zero is $1$.
But the rational function on $\K$ has its actual pole at $\varepsilon$, not
at zero. Indeed its local meromorphic residue at zero is zero. A profile
residue can encode a whole infinitesimal cluster above a standard point;
it is not automatically the local residue at that exact point.

\section{Constructive Weierstrass preparation}

This section supplies the main bridge from complex zero multiplicity to
surcomplex zero geometry. Write the ordinary local variable as $T=z-a$.
All holomorphic coefficient functions in the proof live on a \emph{fixed}
ordinary disk, not on disks whose radii shrink with the exponent.

\begin{theorem}[Preparation for an infinitesimal deformation]\label{thm:weierstrass}
Let $D=D(0,r)$ be an ordinary complex disk. Suppose
\begin{equation}\label{eq:prepinput}
 F(T)=T^d u_0(T)+H(T),\qquad
 u_0\in\OO(D),\quad u_0(T)\neq0\text{ on }D,\quad
 \vprof(H)>0,
\end{equation}
where $d\geq1$. Then there is a unique factorization in $\HH(D)$
\begin{equation}\label{eq:prep}
 F(T)=P(T)U(T)
\end{equation}
with
\[
 P(T)=T^d+p_{d-1}T^{d-1}+\cdots+p_0,\qquad p_j\in\mm,
\]
and $U$ a profile whose exponent-zero coefficient is $u_0$ and whose other
exponents are strictly positive. The lift of $U$ is nowhere zero on $D^\sharp$.
If $S=\supp H$, all positive exponents of $P$ and $U$ may be taken in $M(S)$.
\end{theorem}
\begin{proof}
Set $M=M(S)$, a well-ordered monoid by Neumann's lemma. We seek
\[
 P=T^d+\sum_{\gamma\in M\setminus\{0\}}p_\gamma(T)\tmon^\gamma,
 \qquad
 U=u_0+\sum_{\gamma\in M\setminus\{0\}}u_\gamma(T)\tmon^\gamma,
\]
where every $p_\gamma$ has degree less than $d$, and $u_\gamma$ is holomorphic
on $D$. Let $F_\gamma$ denote the coefficient of $F$, setting it to zero at
exponents not in its support.

At a positive exponent $\gamma$, assume the coefficients at all smaller
exponents have been determined. The known residual is
\begin{equation}\label{eq:prepresidual}
 R_\gamma=F_\gamma-
 \sum_{\substack{\alpha+\beta=\gamma\\\alpha,\beta\in M_{>0}}}
           p_\alpha u_\beta.
\end{equation}
This sum is finite, and every index in it is smaller than $\gamma$.
The coefficient equation to solve is
\begin{equation}\label{eq:preplinear}
 u_0p_\gamma+T^du_\gamma=R_\gamma.
\end{equation}
For a holomorphic function $q$, let $J_{<d}q$ be its Taylor polynomial at zero
through degree $d-1$. The unique solution is
\begin{align}
 p_\gamma&=J_{<d}(R_\gamma/u_0),\label{eq:prepp}\\
 u_\gamma&=
 u_0\,\frac{R_\gamma/u_0-p_\gamma}{T^d}.\label{eq:prepu}
\end{align}
Division by $u_0$ is holomorphic on the same disk. The numerator in
\eqref{eq:prepu} vanishes to order at least $d$ at zero, so the quotient
extends holomorphically there, also on the same disk.

Well-ordered recursion defines all coefficients. Their supports are subsets
of $M$, and all coefficient equations are finite, so these are legitimate
profiles and their product is $F$. For each $j<d$, the coefficients of $T^j$
in the $p_\gamma$ form a Hahn series with positive support, giving $p_j\in\mm$.
The lift of $U$ has nonzero standard part $u_0(\st T)$, so it never vanishes.

For uniqueness, take two factorizations of the stated form. The union of
their positive supports is a well-ordered set. If they differ, choose the
least exponent where a coefficient differs. All lower product terms cancel,
and their differences satisfy
$u_0\Delta p+T^d\Delta u=0$. Dividing by $u_0$ and taking the Taylor polynomial
of degree less than $d$ gives $\Delta p=0$, then $\Delta u=0$, a contradiction.
\end{proof}

The argument remains valid when $S$ is uncountable, or when its positive
scales are not contained in a rank-one subgroup. Even at a limit position
in the ordered monoid, the coefficient equation is a finite equation.
There is no hidden passage to the limit of a sequence of approximants.

\begin{corollary}[The infinitesimal zero cluster]\label{cor:zerocluster}
Under \cref{thm:weierstrass}, $\hhat F$ has exactly $d$ zeros in $D^\sharp$,
counted with analytic multiplicity. All of them belong to $\mm$.
\end{corollary}
\begin{proof}
Algebraic closedness factors $P$ into $d$ linear factors over $\K$, counted
with multiplicity. Every root $\rho$ is infinitesimal. Indeed, if
$\lambda=v(\rho)<0$, then
\[
 v(\rho^d)=d\lambda<j\lambda+v(p_j)=v(p_j\rho^j)
 \quad(j<d),
\]
so the leading term could not cancel. If $\lambda=0$, the leading term again
has valuation zero while all other terms have positive valuation. Both cases
are impossible. A zero root is included in $\mm$ by convention.

Thus all $d$ roots lie in $D^\sharp$. Conversely, the factor $\hhat U$ is
nowhere zero there, so there are no additional zeros. Its germ at any root
is a unit, and multiplication by a unit does not change the order of
vanishing in $\K[[T]]$. Hence polynomial and analytic multiplicities agree.
\end{proof}

\begin{theorem}[Specialization of the zero divisor]\label{thm:divisor}
Let $U$ be connected and $0\neq F\in\HH(U)$ have leading exponent
$\gamma_0$ and leading complex coefficient $f_{\gamma_0}$. For each $a\in U$,
\begin{equation}\label{eq:divisor}
 \sum_{\substack{z\in U^\sharp,\ \hhat F(z)=0\\\st z=a}}
       \ord_z(\hhat F)
       =\ord_a(f_{\gamma_0}).
\end{equation}
The fiber on the left is finite. In divisor notation,
$\st_*\divisor(\hhat F)=\divisor(f_{\gamma_0})$.
\end{theorem}
\begin{proof}
Normalize by $\tmon^{-\gamma_0}$. If the leading coefficient is nonzero at
$a$, \cref{cor:reduction} rules out zeros in its monad. Otherwise its ordinary
holomorphic zero has some finite multiplicity $d$. Choose a disk about $a$
with no other zero of the leading coefficient. On that disk write
$f_{\gamma_0}(a+T)=T^du_0(T)$ with $u_0$ nowhere zero.
Apply \cref{thm:weierstrass,cor:zerocluster}. This yields exactly the stated
finite fiber and its total multiplicity.
\end{proof}

This is a conservation theorem for zero multiplicity, not a statement that
each classical zero has a unique lift. A multiple zero can split into roots
at several infinitesimal separations. If the classical zero is simple,
its lift is unique and is simple.

\section{Argument principle, Rouch\'e, and contour moments}

\subsection{The exact argument principle}

Let $F\in\HH(U)$ be nonzero, with leading coefficient $f_0$ after a constant
monomial normalization. Let $\Gamma$ be a positively oriented ordinary Jordan
contour whose closed interior lies in $U$, and assume $f_0$ has no zeros on
$\Gamma$. Near the contour, write
\[
 F=f_0(1+Q),\qquad\vprof(Q)>0.
\]
The inverse of $F$ is a well-defined profile there, and
\begin{equation}\label{eq:logder}
 \frac{F'}F=\frac{f_0'}{f_0}+
              \bigl(\log_H(1+Q)\bigr)'.
\end{equation}

\begin{theorem}[Surcomplex argument principle]\label{thm:argument}
With the preceding hypotheses,
\begin{equation}\label{eq:argument}
 \frac1{2\pi i}\hoint_\Gamma\frac{F'}F\,dz
  =\sum_{\substack{z\in(\operatorname{int}\Gamma)^\sharp\\\hhat F(z)=0}}
           \ord_z(\hhat F)
  =\sum_{a\in\operatorname{int}\Gamma}\ord_a(f_0).
\end{equation}
In particular, this Hahn-valued integral is exactly a nonnegative ordinary
integer: all its nonconstant Hahn coefficients vanish.
\end{theorem}
\begin{proof}
The derivative of every coefficient of $\log_H(1+Q)$ integrates to zero
around the closed path. This requires no classical logarithm of $f_0$.
The integral in \eqref{eq:logder} is therefore the classical argument-principle
integral of $f_0$, which is the sum of its zero multiplicities inside.
There are finitely many such zeros by compactness and the absence of boundary
zeros. The equality with the actual surcomplex zero count follows from
\cref{thm:divisor}.
\end{proof}

For a quotient $F/G$ of such profiles with leading coefficients nonzero on
the contour, subtracting the two formulas gives the number of zeros minus
the number of poles, with common factors canceled in the local germ sense.

\begin{corollary}[Infinitesimal Rouch\'e theorem]\label{cor:rouche}
If $F\in\HH(U)$ is nonzero and
$\vprof(H)>\vprof(F)$, then $\hhat F$ and $\hhat{F+H}$ have the same total
zero multiplicity in every standard monad in $U^\sharp$. They also have the
same total zero count inside any contour to which \cref{thm:argument} applies.
More generally, after a common normalization to nonnegative exponents, if the
exponent-zero complex coefficients satisfy $|h_0|<|f_0|$ on $\Gamma$, the zero counts of
$F$ and $F+H$ inside the thickened interior are equal.
\end{corollary}
\begin{proof}
In the strictly higher-valuation case the leading complex coefficient is
unchanged, so apply \cref{thm:divisor}. In the second case classical Rouch\'e
applied to $f_0$ and $h_0$ gives equal classical counts and guarantees no
boundary zeros for $f_0+h_0$. Apply \cref{thm:divisor} to both profiles.
\end{proof}

\subsection{Recovering infinitesimal root positions}

Counting is only the zeroth contour moment. Higher moments detect positions
in the surcomplex field, including scales invisible to standard part.

\begin{theorem}[Contour moments of a zero cluster]\label{thm:moments}
In \cref{thm:weierstrass}, let $\Gamma$ be a positively oriented ordinary
circle $|T|=r_1$ with $0<r_1<r$. Let the roots of the prepared polynomial be
$\rho_1,\ldots,\rho_d$, listed with multiplicity. Then for every $k\geq0$,
\begin{equation}\label{eq:moments}
 s_k:=\frac1{2\pi i}\hoint_\Gamma
         T^k\frac{F'(T)}{F(T)}\,dT
      =\sum_{j=1}^d\rho_j^k,
\end{equation}
where $s_0=d$. The moments $s_1,\ldots,s_d$ reconstruct $P$ uniquely.
\end{theorem}
\begin{proof}
Differentiate $F=PU$ to get $F'/F=P'/P+U'/U$. The latter term is a
holomorphic profile on $D$, because $U$ has nowhere-zero leading coefficient.
Thus its product with $T^k$ integrates to zero.

In $\K(T)$, logarithmic differentiation of the factorization of $P$ gives
$P'/P=\sum_j(T-\rho_j)^{-1}$. On the ordinary circle, each $\rho_j$ is
infinitesimal, so
\[
 \frac1{T-\rho_j}=\sum_{n\geq0}\rho_j^nT^{-n-1}
\]
is a well-defined profile. The coefficientwise contour integral of
$T^k/(T-\rho_j)$ is $2\pi i\rho_j^k$, proving \eqref{eq:moments}.

Let $e_0=1$ and define the elementary symmetric functions recursively by
Newton's identities
\begin{equation}\label{eq:newtonidentities}
 k e_k=\sum_{j=1}^k(-1)^{j-1}e_{k-j}s_j,\qquad 1\leq k\leq d.
\end{equation}
Characteristic zero permits the divisions by $k$, and
$P(T)=T^d-e_1T^{d-1}+e_2T^{d-2}-\cdots+(-1)^de_d$.
\end{proof}

Thus the prepared polynomial has two constructive descriptions: a recursion
using Taylor division, and a finite collection of Hahn-valued contour moments.
Neither description discards infinitesimal corrections to the roots.

\section{Biholomorphisms and a thickened Riemann mapping theorem}

\begin{theorem}[Lifting an infinitesimal deformation of a biholomorphism]
\label{thm:biholo}
Let $f_0:U\to V$ be an ordinary biholomorphism of complex domains, and let
$F=f_0+H\in\HH(U)$ with $\vprof(H)>0$. There is a profile
$G=g_0+J\in\HH(V)$, where $g_0=f_0^{-1}$ and $\vprof(J)>0$, such that
\[
 F\circ G=\id_V,\qquad G\circ F=\id_U.
\]
Consequently, $\hhat F:U^\sharp\to V^\sharp$ is a bijection with germ-analytic
inverse $\hhat G$. Its derivative never vanishes.
\end{theorem}
\begin{proof}
Let $M$ be the positive monoid generated by the support of $H$. Seek
$J=\sum_{\gamma\in M_{>0}}j_\gamma\tmon^\gamma$.
In the coefficient equation at exponent $\gamma$ in $F(g_0+J)=\id$, the
only term involving the new unknown $j_\gamma$ is
$f_0'(g_0)j_\gamma$. All other terms involve lower exponents or known
coefficients of $H$. Neumann's lemma makes each such equation finite.
Because $f_0'(g_0)$ is nowhere zero on $V$, it determines a holomorphic
$j_\gamma$ on that same domain. Well-ordered recursion constructs $G$,
and \cref{thm:profilecomposition} justifies the composition identity.

The standard part of $\hhat F(a+\eta)$ is $f_0(a)$, so its image is in
$V^\sharp$. The evaluated identity $\hhat F\circ\hhat G=\id$ gives
surjectivity. For $w=b+\varepsilon\in V^\sharp$, the profile $F-w$ has leading
complex coefficient $f_0-b$, with a unique simple zero at $a=g_0(b)$.
By \cref{thm:divisor}, there is exactly one solution of $\hhat F(z)=w$ in
that monad, and there are no solutions in other monads. Thus $\hhat F$ is
injective as well. Evaluation is faithful, so the other profile composition
is the identity. Finally $\st\bigl(\hhat{F'}(a+\eta)\bigr)=f_0'(a)\neq0$.
\end{proof}

\begin{corollary}[Riemann mapping for infinitesimal thickenings]
\label{cor:riemann}
If $U\subsetneq\C$ is nonempty and simply connected, then $U^\sharp$ is
biholomorphic in the profile sense to $\D^\sharp$. Every ordinary Riemann map
lifts, and every infinitesimal holomorphic-profile deformation of that map
is again such a biholomorphism onto the same thickened target.
\end{corollary}
\begin{proof}
Use the classical Riemann mapping theorem to obtain a biholomorphism
$f_0:\D\to U$, and apply \cref{thm:biholo}.
\end{proof}

This is a theorem about domains of the form $U^\sharp$, not a Riemann mapping
theorem for arbitrary sharp-open subsets of $\K$. In particular,
\[
 \D^\sharp=\{z\in\OO:|\st z|<1\}
 \quad\text{is strictly smaller than}\quad\{z\in\K:|z|<1\}.
\]
The latter set also contains points such as $1-\varepsilon$ with positive
infinitesimal $\varepsilon$, whose standard part lies on the unit circle.
None of that boundary monad belongs to $\D^\sharp$.

\subsection{Schwarz's lemma and the importance of the category}

\begin{proposition}[Schwarz's lemma for canonical classical lifts]
\label{prop:schwarz}
Let $f:\D\to\D$ be ordinary holomorphic with $f(0)=0$, and let $\hhat f$
be its canonical lift, with no additional Hahn coefficients. Then
\[
 |\hhat f(z)|\leq|z|\qquad(z\in\D^\sharp).
\]
For nonzero $z$, equality is possible only when $f(z)=e^{i\theta}z$ for an
ordinary real $\theta$, in which case equality holds everywhere.
\end{proposition}
\begin{proof}
Classical Schwarz's lemma writes $f(z)=zg(z)$, with $g$ holomorphic on $\D$.
Unless $f$ is a rotation, the classical maximum principle gives $|g(a)|<1$
for every $a\in\D$, including $a=0$. Therefore
$|\hhat g(a+\eta)|<1$: its standard modulus is strictly below one by an
ordinary positive margin. Lifting the product identity proves the strict
inequality at every nonzero surcomplex point. A rotation lifts exactly.
\end{proof}

The same conclusion fails for general profiles, even very simple ones.
For $t=\tmon>0$, the profile
\[
 F(z)=(1+t)z
\]
is a biholomorphism of $\D^\sharp$ onto itself, fixes zero, and satisfies
$|\hhat F(z)|>|z|$ for every nonzero $z$. Its leading complex map is the
identity, so \cref{thm:biholo} applies. Likewise
\[
 G(z)=z+t z^2
\]
is a nonidentity biholomorphism of $\D^\sharp$ fixing zero and satisfying
$G'(0)=1$. Thus even the usual normalization that would characterize an
ordinary disk automorphism does not eliminate infinitesimal profile
freedom. There is no contradiction with \cref{prop:schwarz}, whose hypothesis
specifies the canonical lift of an ordinary disk map.

\section{Newton geometry: seeing zeros at infinite as well as finite scales}

Profiles over $\C$ naturally describe finite arguments. Arbitrary centers and
scales are reached through affine charts $z=A+R w$, with $A\in\K$ and
$R\in\K^\times$. The local theory is unchanged by such a substitution.
For polynomials there is a particularly efficient global scale detector.

\begin{definition}
For a nonzero polynomial $P(Z)=\sum_{j=0}^d a_j Z^j$ and $\lambda\in\No$, put
\[
 g_P(\lambda)=\min_{a_j\neq0}\{v(a_j)+j\lambda\}.
\]
This is the weighted Gauss valuation of $P$ at scale $\lambda$.
\end{definition}

\begin{theorem}[Newton polygon in the surcomplex field]\label{thm:newton}
For nonzero polynomials $P,Q$,
\[
 g_{PQ}(\lambda)=g_P(\lambda)+g_Q(\lambda).
\]
If
$P(Z)=a_d\prod_{j=1}^d(Z-\rho_j)$, with roots listed with multiplicity, then
\begin{equation}\label{eq:tropical}
 g_P(\lambda)=v(a_d)+\sum_{j=1}^d\min\{\lambda,v(\rho_j)\},
\end{equation}
where $v(0)=+\infty$. Thus its changes of slope occur at the valuations of
the nonzero roots, with drops equal to their multiplicities. Equivalently,
the lower Newton-polygon edges through the finite points $(j,v(a_j))$,
with $a_j\neq0$, have slopes $-v(\rho_j)$ for nonzero roots and horizontal
lengths equal to the corresponding multiplicities. Zero roots are recorded
by the largest power of $Z$ dividing $P$.
\end{theorem}
\begin{proof}
Replace $Z$ by $\tmon^\lambda W$ and divide $P$ by
$\tmon^{g_P(\lambda)}$. All coefficients of the resulting polynomial lie in
$\OO$, and its reduction in $\C[W]$ is nonzero. Do the same for $Q$.
The product of their nonzero reductions is nonzero, because $\C[W]$ is an
integral domain. The normalized product therefore has coefficient valuation
zero, which proves multiplicativity. Applying it to the linear factors gives
\eqref{eq:tropical}.

Each summand $\min\{\lambda,v(\rho_j)\}$ has slope one before its breakpoint
and zero afterwards. This proves the multiplicity assertion. Taking the
lower envelope of the finitely many affine functions $v(a_j)+j\lambda$
is exactly the dual description of the lower Newton polygon. Only finite
arithmetic and division by positive integers in the ordered group are used.
\end{proof}

A direct necessary condition for a nonzero root with valuation $\lambda$ is
that the minimum defining $g_P(\lambda)$ occur at least twice; otherwise the
unique least-valued term could not cancel. The factorization proof shows how
this familiar tropical condition becomes an exact multiplicity statement.

\begin{example}[Two infinitesimal ranks in one cubic]\label{ex:cubic}
Consider
\[
 P(Z)=Z^3-\tmon^2 Z-\tmon^{\omega+2}.
\]
Its lower Newton polygon has edges joining
$(0,\omega+2)$ to $(1,2)$ and $(1,2)$ to $(3,0)$.
There is one root of valuation $\omega$ and two roots of valuation $1$.
Set $Z=\tmon Y$ and $q=\tmon^{\omega-1}$. Then
\[
 P(\tmon Y)=\tmon^3(Y^3-Y-q).
\]
The three simple ordinary roots $0,1,-1$ of $Y^3-Y$ lift uniquely. Their first
terms are
\begin{align*}
 Y_0&=-q-q^3-3q^5+O(q^7),\\
 Y_+&=1+\tfrac12q-\tfrac38q^2+\tfrac12q^3+O(q^4),\\
 Y_-&=-1+\tfrac12q+\tfrac38q^2+\tfrac12q^3+O(q^4).
\end{align*}
Multiplying by $\tmon$ displays the finer scales of the original roots:
\[
 Z_0=-\tmon^\omega-\tmon^{3\omega-2}+\cdots,\qquad
 Z_\pm=\pm\tmon+\tfrac12\tmon^\omega
          \mp\tfrac38\tmon^{2\omega-1}+\cdots.
\]
Here $q$ is a formal parameter subsequently evaluated as an infinitesimal;
$O(q^N)$ denotes a formal series divisible by $q^N$, not a claim of sharp
sequential convergence. The gap between valuations $1$ and $\omega$ is
not visible in a rank-one Archimedean value group.
\end{example}

As a simpler domain warning, $1+\tmon z$ has no zero in $\C^\sharp=\OO$,
although as a polynomial on $\K$ it has the root $-\tmon^{-1}$.
The profile theorem counts all zeros in its specified thickening, not zeros
outside that domain. A rescaling reveals the infinite root immediately.

\section{Laurent expansions, removability, and growth}

\subsection{Two different meanings of a Laurent expansion}

For $F=\sum f_\gamma\tmon^\gamma\in\HH(A)$ on an ordinary annulus
$A=\{r<|z-a|<R\}$, every coefficient has its ordinary Laurent expansion
\[
 f_\gamma(z)=\sum_{n\in\Z}c_{\gamma,n}(z-a)^n,
 \qquad
 c_{\gamma,n}=\frac1{2\pi i}\int_{|\zeta-a|=s}
       \frac{f_\gamma(\zeta)}{(\zeta-a)^{n+1}}\,d\zeta,
\]
for $r<s<R$. This is a \emph{coefficientwise Laurent expansion}.
The inner Laurent series is first evaluated in the usual complex sense,
and the resulting holomorphic functions are then lifted by
\eqref{eq:lift}. There is no general justification for interpreting the
whole array $(\gamma,n)$ as one strong Hahn sum at a nonzero ordinary
complex argument. Infinitely many ordinary scalar terms would contribute
at exponent $\gamma$.

By contrast, a \emph{meromorphic germ} at an arbitrary $a\in\K$ is an element
of $\K((T))$, with $T=z-a$ and only finitely many negative powers. Its positive
part is evaluated by \cref{thm:scale}, and its finite principal part by field
arithmetic. This gives a local meromorphic function on a sufficiently small
punctured sharp neighborhood. The formal residue calculus for these germs
is developed in \cref{sec:formalresidue}.

\begin{proposition}[Coefficientwise removable singularities]\label{prop:removal}
Suppose $F=\sum f_\gamma\tmon^\gamma$ has common support and is holomorphic
on an ordinary punctured disk $D\setminus\{a\}$. It extends to a profile in
$\HH(D)$ if and only if each coefficient $f_\gamma$ has a removable singularity
at $a$. In particular, extension holds if every $f_\gamma$ is bounded near $a$,
with a bound allowed to depend on $\gamma$.
\end{proposition}
\begin{proof}
Any profile extension restricts coefficientwise to the given functions,
by normal-form uniqueness at ordinary points, so necessity is immediate.
Conversely, the ordinary holomorphic extensions have support contained in
the same well-ordered set and therefore form a profile on $D$.
For bounded coefficients the classical removable singularity theorem applies;
alternatively its Cauchy estimate makes every negative Laurent coefficient
zero by letting the ordinary contour radius tend to zero.
\end{proof}

Uniformly bounded pole order has a stronger consequence. If for one finite
integer $d$ all functions $(z-a)^d f_\gamma(z)$ extend holomorphically across
$a$, then $(z-a)^dF$ is a holomorphic profile on $D$. Its lift divided by
$(z-a)^d$ defines a meromorphic function on $D^\sharp\setminus\{a\}$ with a
pole of order at most $d$ at the exact point $a$. In this case the profile
residue agrees with the local meromorphic residue there. The common finite
pole-order condition excludes the displaced-pole phenomenon in
\eqref{eq:movingpole}.

\subsection{What a valid Liouville theorem says}

\begin{theorem}[Coefficientwise polynomial growth]\label{thm:growth}
Let $F=\sum f_\gamma\tmon^\gamma\in\HH(\C)$, and fix an ordinary integer
$d\geq0$. Suppose that for each $\gamma$ there is a real constant $C_\gamma$
such that
\[
 |f_\gamma(z)|\leq C_\gamma(1+|z|)^d\qquad(z\in\C).
\]
Then $F$ is a polynomial of degree at most $d$ with coefficients in $\K$.
In particular, if every coefficient is bounded on $\C$, then $F$ is constant.
\end{theorem}
\begin{proof}
For fixed $a\in\C$ and $n>d$, the ordinary Cauchy estimate on a circle of
radius $R$ gives
\[
 |f_\gamma^{(n)}(a)|\leq
 \frac{n!C_\gamma(1+|a|+R)^d}{R^n}.
\]
Let $R$ tend to infinity in the ordinary real sense. The right side tends
to zero, so each $f_\gamma$ is a complex polynomial of degree at most $d$.
Write $f_\gamma(z)=\sum_{j=0}^d c_{\gamma,j}z^j$. For every $j$, the series
$A_j=\sum_\gamma c_{\gamma,j}\tmon^\gamma$ has support contained in the original
support and defines an element of $\K$. Regrouping the finite sum over $j$
gives $F(z)=\sum_{j=0}^d A_jz^j$.
\end{proof}

A pointwise surreal bound is much weaker. The profile $F(z)=\tmon z$ is
nonconstant on $\C^\sharp=\OO$, yet $|\hhat F(z)|<1$ for every finite $z$.
Its coefficient $z$ is unbounded on $\C$; the infinitesimal factor hides that
unboundedness from any single ordinary positive pointwise bound.

Likewise,
\[
 F(z)=\tmon e^{1/z}
\]
is bounded in modulus by $1$ on $(D\setminus\{0\})^\sharp$, because at every
point of that domain its exponential factor is finite. It has no holomorphic
profile extension over zero, since its coefficient $e^{1/z}$ has an essential
singularity. Notice that $(D\setminus\{0\})^\sharp$ deletes the entire zero
monad, not merely the point zero. The two punctured domains must not be
identified.

Even on a full sharp punctured neighborhood, boundedness alone is insufficient
for arbitrary germ-analytic functions. For $z\neq0$, define $b(z)$ to be $1$
when the surreal $v(z)$ is an ordinal, and $0$ otherwise. Valuation is locally
constant away from zero, so $b$ is locally constant and hence germ-analytic.
Ordinals are cofinal in $\No$, and both $\tmon^\alpha$ and
$\tmon^{\alpha+1/2}$ can be made arbitrarily small by choosing a sufficiently
large ordinal $\alpha$. Their $b$-values differ. Thus $b$ is bounded and has
no sharp limit at zero. This is an explicit failure of an unrestricted
Riemann removable-singularity theorem.

\section{Exponentials and logarithms: local strength, global freedom}

\subsection{The canonical exponential on a large strip}

We use the standard Kruskal--Gonshor exponential
$\exp_{\No}:\No\to\No_{>0}$ and its inverse $\log_{\No}$. It is an increasing
group isomorphism extending the real exponential, and on infinitesimals it
agrees with its Hahn power series \cite{Gonshor,BMder}.
For a finite surcomplex number $a+\eta$ with $a\in\C$ and $\eta\in\mm$, define
\[
 \exp_{\fin}(a+\eta)=e^a\exp_H(\eta).
\]
The formal addition law proves that this is an exponential homomorphism on
the additive group $\OO$.

Consider the open additive class
\[
 \mathcal A=\No+i(\No\cap\OO).
\]
Its elements may have an arbitrary, even infinite, real part, but have a
finite imaginary part. Define
\begin{equation}\label{eq:strip-exp}
 E(x+iy)=\exp_{\No}(x)\,\exp_{\fin}(iy),\qquad x+iy\in\mathcal A.
\end{equation}

\begin{theorem}[An exponential covering all nonzero values]\label{thm:strip-exp}
The map $E:\mathcal A\to\K^\times$ is a surjective group homomorphism,
is germ-analytic, and satisfies
\[
 E'=E,\qquad |E(x+iy)|=\exp_{\No}(x),\qquad
 \ker E=2\pi i\Z.
\]
Thus there is an exact sequence of additive and multiplicative class groups
\[
 0\longrightarrow2\pi i\Z\longrightarrow\mathcal A
       \xrightarrow{\ E\ }\K^\times\longrightarrow1.
\]
\end{theorem}
\begin{proof}
The addition law follows separately from the surreal real exponential and
the finite complex exponential. If $h\in\mm$, its real and imaginary parts
are infinitesimal, and the formal exponential product identity gives
\begin{equation}\label{eq:exp-local}
 E(z+h)=E(z)\exp_H(h).
\end{equation}
This proves germ-analyticity and $E'=E$.
Conjugation and the addition law show that $\exp_{\fin}(iy)$ has modulus one,
which gives the modulus formula. If $E(x+iy)=1$, the injectivity of the real
exponential forces $x=0$. Write $y=y_0+\varepsilon$ with $y_0\in\R$.
Taking standard part gives $e^{iy_0}=1$, hence $y_0\in2\pi\Z$; then
$\exp_H(i\varepsilon)=1$ forces $\varepsilon=0$ by the inverse logarithm.
This proves the kernel assertion.

For surjectivity, write a nonzero $w$ in the unique form
$w=c\tmon^\gamma(1+\eta)$ from \eqref{eq:leadingunit}. Choose any ordinary
complex logarithm $\ell(c)$ of $c$. Then
\begin{equation}\label{eq:logsection}
 L(w)=\log_{\No}(\tmon^\gamma)+\ell(c)+\log_H(1+\eta)
\end{equation}
belongs to $\mathcal A$ and satisfies $E(L(w))=w$. Its imaginary part is finite,
because the first term is real and the other two have finite imaginary part.
\end{proof}

This surjectivity is worth emphasizing: arbitrary nonzero surcomplex values
do not require an infinite imaginary logarithm. An arbitrary real logarithmic
magnitude together with a finite phase is sufficient.

\subsection{A global locally analytic logarithm}

Choose once and for all a single-valued set-theoretic logarithm
$\ell:\C^\times\to\C$, for example a principal-value convention, and define
$L$ by \eqref{eq:logsection} for every $w\in\K^\times$.

\begin{theorem}[Global logarithmic section]\label{thm:globallog}
The function $L:\K^\times\to\mathcal A$ is germ-analytic on its entire domain,
with
\[
 E(L(w))=w,\qquad L'(w)=\frac1w.
\]
In particular, $1/w$ has a global primitive in the category of germ-analytic
functions on $\K^\times$.
\end{theorem}
\begin{proof}
The exponential identity was already proved. Fix $w\neq0$ and take
$h$ with $v(h)>v(w)$. Then $h/w\in\mm$, and $w+h$ has the same leading
exponent and the same leading complex coefficient as $w$. In
\eqref{eq:logsection}, the first two terms therefore do not change. Its unit
part changes from $1+\eta$ to $(1+\eta)(1+h/w)$. The formal logarithm's
addition law gives the exact local identity
\begin{equation}\label{eq:global-log-local}
 L(w+h)-L(w)=\log_H(1+h/w).
\end{equation}
This is a power series in $h$ with linear coefficient $1/w$.
\end{proof}

There is no conflict with
$\hoint_{|z|=1}dz/z=2\pi i$. The function $L$ is not a holomorphic profile on
$\C^\times$. At ordinary complex points it equals the chosen branch
$\ell$, which cannot be holomorphic on all of $\C^\times$. The coefficientwise
fundamental theorem applies to profile primitives, not to every sharp-local
primitive. The usual monodromy obstruction survives in the profile category
and disappears in the more permissive local category.

\subsection{Many globally germ-analytic exponentials}

The normal form splits every real surreal uniquely as
\[
 y=y^{\infp}+y^{\fin},
\]
where $y^{\infp}$ consists of all strictly infinite monomials and
$y^{\fin}\in\No\cap\OO$ consists of the remaining terms. Both projections
are additive. Let $c_\omega(y)\in\R$ be the coefficient of the Conway monomial
$\omega$ in $y$. For any ordinary real $\lambda$, put
\[
 P_\lambda(y)=y^{\fin}+\lambda c_\omega(y),\qquad
 \mathcal E_\lambda(x+iy)=
       \exp_{\No}(x)\exp_{\fin}\bigl(iP_\lambda(y)\bigr).
\]

\begin{proposition}[Nonuniqueness even with the exponential laws]
\label{prop:manyexp}
Each $\mathcal E_\lambda$ is a germ-analytic function on all of $\K$ satisfying
\[
 \mathcal E_\lambda(z+w)=\mathcal E_\lambda(z)\mathcal E_\lambda(w),\qquad
 \mathcal E_\lambda'=\mathcal E_\lambda,\qquad
 \mathcal E_\lambda(0)=1.
\]
It extends the ordinary complex exponential and the surreal real exponential,
commutes with conjugation, and has
$|\mathcal E_\lambda(z)|=\exp_{\No}(\Ree z)$.
Nevertheless,
\[
 \mathcal E_\lambda(i\omega)=e^{i\lambda}.
\]
Thus these conditions do not select a unique global surcomplex exponential.
\end{proposition}
\begin{proof}
Additivity of $P_\lambda$ gives the group law, and it is the identity on finite
real surreals, giving both extension properties. Conjugation and modulus
follow as in \cref{thm:strip-exp}. For an infinitesimal real increment
$\varepsilon$, one has $P_\lambda(y+\varepsilon)=P_\lambda(y)+\varepsilon$.
Consequently \eqref{eq:exp-local} holds for $\mathcal E_\lambda$ at every point
of $\K$, proving its analytic power-series representation and derivative.
Finally $P_\lambda(\omega)=\lambda$.
\end{proof}

For $\lambda=0$, the kernel is
$i(\mathbb J+2\pi\Z)$, where $\mathbb J$ is the additive class of purely
infinite real surreals. These constructions are legitimate but choice-dependent
global analytic structures. They are not a claim that infinite imaginary
phases have a uniquely preferred extension. Local differential equations
cannot determine data on separated sharp regions.

\section{Worked deformations and explicit calculations}

The expansions below are exact formal power-series identities followed by
Hahn evaluation at an infinitesimal parameter. They are not numerical
approximations to real or complex solutions. Truncations only abbreviate the
higher formal coefficients.

\subsection{Motion of a simple zero}

Let $f_0,g,h$ be ordinary holomorphic near $a$, with $f_0(a)=0$ and
$A=f_0'(a)\neq0$. For a positive-valuation parameter $q$, consider
\[
 F(z)=f_0(z)+qg(z)+q^2h(z).
\]
Its unique zero in the monad of $a$ has the expansion
\begin{equation}\label{eq:simplemotion}
 \rho=a-\frac{g(a)}Aq+
 \left(
 \frac{g(a)g'(a)}{A^2}
 -\frac{f_0''(a)g(a)^2}{2A^3}
 -\frac{h(a)}A
 \right)q^2+O(q^3).
\end{equation}
Indeed, substituting $\rho=a+c_1q+c_2q^2+\cdots$ yields the equations
\[
 Ac_1+g(a)=0,\qquad
 Ac_2+\tfrac12f_0''(a)c_1^2+g'(a)c_1+h(a)=0.
\]
The higher coefficients are determined successively by the same invertible
linear coefficient $A$. The zero-counting theorem guarantees that this
constructed formal root is the only actual zero in that monad.

\subsection{A double zero of a transcendental profile}
\label{subsec:doubleexample}

Take
\[
 F(T)=T^2+q e^T,\qquad v(q)>0.
\]
The leading complex function has a double zero at zero. The preparation
recursion with $u_0=1$ gives
\begin{align}
 P(T)&=T^2+q(1+T)-q^2\left(\tfrac12+\tfrac23T\right)+O(q^3),
       \label{eq:exampleP}\\
 U(T)&=1+q\,\frac{e^T-1-T}{T^2}+O(q^2).
       \label{eq:exampleU}
\end{align}
For the second coefficient in $P$, the residual is
$-(1+T)(e^T-1-T)/T^2$; its constant and linear terms are $-1/2$ and $-2/3$.
All apparent singularities in $U$ are removable.

Write $q=s^2$ and put $T=sY$. The equation becomes
$Y^2+e^{sY}=0$. At $s=0$ its roots $Y=\pm i$ are simple, so the formal
implicit function theorem constructs both branches. Matching coefficients
gives the two roots
\begin{equation}\label{eq:doubleroots}
 \rho_\pm=\pm is-\tfrac12s^2\mp\tfrac{3i}{8}s^3
                 +\tfrac13s^4+O(s^5).
\end{equation}
The prepared polynomial therefore records, to this order,
\[
 \rho_++\rho_-=-q+\tfrac23q^2+O(q^3),\qquad
 \rho_+\rho_-=q-\tfrac12q^2+O(q^3).
\]
These agree with \eqref{eq:exampleP}. For any fixed ordinary circle around
zero, \cref{thm:moments} gives
\[
 \frac1{2\pi i}\hoint \frac{F'}F\,dT=2,\qquad
 \frac1{2\pi i}\hoint T\frac{F'}F\,dT=-q+\tfrac23q^2+O(q^3),
\]
and
\[
 \frac1{2\pi i}\hoint T^2\frac{F'}F\,dT=-2q+2q^2+O(q^3).
\]
The zeroth moment is an exact integer; the higher moments retain the
infinitesimal displacement information.

\subsection{Arbitrary coefficient growth at infinitesimal arguments}

The ordinary divergent series
\[
 A(T)=\sum_{n\geq0}n!\,T^n
\]
is well-defined as a Hahn-analytic function at every infinitesimal $T$:
its coefficients are complex scalars, and \cref{cor:infseries} applies.
As a formal identity,
\begin{equation}\label{eq:factorialODE}
 T^2 A'(T)+(T-1)A(T)+1=0.
\end{equation}
For $n\geq1$, the coefficient equation is
$(n-1)(n-1)!+(n-1)!-n!=0$; the constant term also cancels.
Thus the divergent classical formal solution becomes an actual surcomplex
local analytic solution, without introducing Borel summation or identifying
Hahn summation with an ordinary real limiting process.
It does not thereby become a holomorphic profile over any ordinary
neighborhood of zero. Indeed, matching all derivatives at zero would force
every nonzero-exponent coefficient of such a profile to have zero Taylor
series and hence to vanish identically. Its exponent-zero coefficient would
then be an ordinary holomorphic function with Taylor coefficients $n!$,
which is impossible by the ordinary Cauchy estimates.

\section{A compatible second derivation}
\label{sec:derivations}

Surreal transseries calculus offers a second, different derivative. Let
$\partial:\K\to\K$ be a field derivation killing $\C$ and \emph{strongly
additive}: whenever a family is Hahn summable, the family of its derivatives
is summable and the derivative of the sum is the sum of the derivatives.
The Berarducci--Mantova derivation on $\No$, extended with $\partial i=0$,
is an example \cite{BMder}.

This derivation can change a constant such as $\omega$. The variable derivative
$D_z$ used above cannot: $D_z\omega=0$. The two derivatives nevertheless admit
a precise compatibility law.

\begin{theorem}[Coefficient and variable differentiation]\label{thm:twoderiv}
For $F=\sum f_\gamma\tmon^\gamma\in\HH(U)$, define
\[
 \partial_cF=\sum_\gamma f_\gamma\,\partial(\tmon^\gamma),
\]
expanding the derivatives into normal form and regrouping. This is a
well-defined holomorphic profile, and
\[
 D_z\partial_cF=\partial_cD_zF.
\]
At every $z\in U^\sharp$, its evaluation satisfies the total derivative law
\begin{equation}\label{eq:two-deriv}
 \partial\bigl(\hhat F(z)\bigr)
   =\hhat{\partial_cF}(z)+\hhat{F'}(z)\,\partial z.
\end{equation}
\end{theorem}
\begin{proof}
The monomial family $(\tmon^\gamma)_{\gamma\in S}$ is summable. Strong
additivity says that $(\partial\tmon^\gamma)_{\gamma\in S}$ is summable as
well. Its joint support is therefore well ordered, with only finitely many
contributing $\gamma$ at each new exponent. Multiplication by the complex
coefficient functions leaves that support bound unchanged. Every new
coefficient is a finite complex linear combination of holomorphic functions,
so $\partial_cF$ is a profile. Ordinary variable differentiation commutes
with these finite linear combinations.

Write $z=a+\eta$. Apply strong additivity to the summable family in
\eqref{eq:lift}, and use $\partial a=0$. Leibniz gives
\[
 \partial\bigl(\tmon^\gamma\eta^n\bigr)
   =\partial(\tmon^\gamma)\eta^n
       +\tmon^\gamma n\eta^{n-1}\partial\eta.
\]
The first family regroups to $\hhat{\partial_cF}(z)$: its support is the
well-ordered derivative support plus the positive monoid generated by
$\supp\eta$. The second regroups to $\hhat{F'}(z)\partial\eta$.
All these families are summable by Neumann's lemma and multiplication of a
summable family by a fixed element. Since $\partial\eta=\partial z$, this
is \eqref{eq:two-deriv}.
\end{proof}

\begin{corollary}[Differentiating a moving simple zero]
If $\rho\in U^\sharp$ is a simple zero of $\hhat F$, then
\[
 \partial\rho=-\frac{\hhat{\partial_cF}(\rho)}{\hhat{F'}(\rho)}.
\]
\end{corollary}
\begin{proof}
Differentiate the equality $\hhat F(\rho)=0$ using \eqref{eq:two-deriv},
and divide by the nonzero variable derivative at the simple zero.
\end{proof}

This connects infinitesimal zero geometry to asymptotic differential algebra
without confusing differentiation of the argument with differentiation of
the coefficient field.

\section{Scope, structural lessons, and further development}
\label{sec:scope}

\subsection{Which hypotheses recover which theorems?}

The following table summarizes the actual scope of the results. In particular,
``analytic'' is never used as an unqualified substitute for every classical
hypothesis at once.

\renewcommand{\arraystretch}{1.18}
\begin{longtable}{@{}>{\raggedright\arraybackslash}p{0.27\textwidth}p{0.67\textwidth}@{}}
\toprule
Classical theme & Surcomplex statement and essential convention\\
\midrule
\endfirsthead
\toprule
Classical theme & Surcomplex statement and essential convention\\
\midrule
\endhead
Power-series calculus & Every element of $\K[[T]]$ has a local Hahn-summable
realization. Its radius may have to be smaller than every member of a
specified set of positive scales. Partial sums need not converge sharply.\\
Inverse, implicit, and ramification theorems & These hold for germ-analytic
functions at arbitrary finite or infinite points. They are consequences of
formal algebra together with the scale-separation theorem.\\
Identity theorem & A holomorphic profile over a connected ordinary domain is
determined by one full jet, or by its values on any nonempty sharp-open set.
An arbitrary globally germ-analytic function has no such uniqueness.\\
Cauchy and residue formulas & Contours are ordinary complex contours and
integration is coefficientwise. Cauchy formulas evaluate the lifted function
at infinitesimally displaced arguments. A profile residue may represent an
entire moving pole cluster.\\
Preparation and zero counting & A positive-valuation deformation of a complex
zero of multiplicity $d$ has a unique monic degree-$d$ prepared polynomial
and exactly $d$ zeros, with multiplicities, in that monad.\\
Open mapping and maximum modulus & They hold locally at every point where the
analytic germ is nonconstant. A nonconstant profile on a connected base has
this property everywhere. The modulus is the surreal Euclidean modulus,
not an ultrametric absolute value.\\
Riemann mapping & Ordinary conformal maps and their positive-valuation profile
deformations give biholomorphisms between infinitesimal thickenings. No claim
is made about all sharp-open subclasses or the missing boundary monads.\\
Schwarz rigidity & The usual sharp inequality and equality case hold for
canonical lifts of ordinary disk maps. Infinitesimal profile deformations can
violate both the inequality and normalized uniqueness.\\
Liouville and removability & Coefficientwise complex bounds give valid
statements. A single surcomplex bound on the evaluated function can hide an
unbounded or singular coefficient and is not sufficient.\\
Logarithms & A global logarithm exists in the very local germ-analytic
category. It is not a holomorphic profile on $\C^\times$, where the usual
contour obstruction survives.\\
\bottomrule
\end{longtable}
\renewcommand{\arraystretch}{1}

\subsection{What uses the full class of surreal numbers?}

Most of the profile arguments use only a Hahn presentation, the ordered
abelian group law on exponents, Neumann's lemma, and ordinary complex analysis.
Preparation itself is a support-controlled algebraic recursion. Algebraic
closedness is additionally used to split its polynomial and identify actual
roots. Consequently these arguments also apply in an algebraically closed
set-sized Hahn field with coefficient field $\C$, whenever the stated
supports and operations belong to that field.

Two assertions are specifically strengthened by working in the full $\No$.
First, a set of coefficient supports, however large, has a surreal upper and
lower bound. This is the reason \emph{every} formal series acquires a local
radius in \cref{thm:scale}. Second, the same bounding property makes every set
subset sharply closed and discrete. These are complementary manifestations
of the enormous range of scales: local algebra gains room, while ordinary
set-indexed convergence loses power. Neither assertion should be transferred
without examination to a fixed set-sized subfield.

The construction does not prove that arbitrary sharply complex-differentiable
functions are germ-analytic. Such a converse would be a separate analogue of
Goursat's theorem, with its own regularity and domain hypotheses. Here the
power-series condition is explicit from the outset, and each global theorem
specifies the additional profile coherence that it needs.

\subsection{Natural extensions of the framework}

There are several concrete directions in which the theory can be enlarged
without changing its basic discipline.

\paragraph{An atlas of scales.}
The affine changes $z=A+Rw$ already transport local theorems to infinite
centers and arbitrary surreal radii. A global atlas can retain the common
support data on overlaps, rather than treating all independently analytic
local functions as one coherent object. This would incorporate the boundary
monads omitted by $U^\sharp$ and organize the multiple scales visible in
Newton polygons.

\paragraph{Locally controlled profiles.}
The present $\HH(U)$ is a presheaf with a global support requirement. For
infinite gluing, one can instead require a common well-ordered support on each
member of a specified ordinary open cover, together with exact compatibility
on overlaps. To preserve integration, every compact contour must admit a
finite subcover with a jointly admissible support. Stating this condition is
essential: arbitrary gluing need not preserve a single Hahn support.

\paragraph{Deformed contours.}
For a profile deformation of a parametrized ordinary contour, the natural
candidate integral is obtained by pullback and coefficientwise integration.
Its validity requires one support bound for the pullback, its parameter
derivative, and the products involved. The fixed-contour Cauchy formulas proved
here supply the reference case; they do not silently authorize integration
along every sharp-continuous class-parametrized path.

\paragraph{Asymptotic differential geometry.}
The two-derivation identity \eqref{eq:two-deriv} gives a usable starting point
for differential equations whose coefficients and zeros both move through
surreal scales. The coefficientwise preparation recursion then separates the
finite algebra of a zero cluster from the asymptotic evolution of its
coefficients.

\subsection{Reproducibility and the status of the constructions}

The companion script \texttt{verify\_examples.py} checks the finite symbolic
expansions in the worked examples: the cubic root branches, the double-zero
preparation and root splitting, the contour moments, the quadratic inverse,
and the factorial-series differential identity. These are exact symbolic
checks, not numerical tests at a large finite surrogate for $\omega$.
They verify the printed finite calculations; they are not a formal proof
verification of the general theorems.

The foundational surreal facts, Hahn summation lemma, classical complex
results, and existing surreal derivation are attributed to their sources.
The analytic structures and deductions are developed explicitly here, with
proofs. This is a proposed framework and research exposition, not a claim
that every formulation has no antecedent in the literature, nor that one
unique canonical global theory has been selected. Its central mathematical
point is positive: support control and ordinary holomorphic coherence recover
strong complex-analytic structure without suppressing genuinely surreal
infinitesimal geometry.

\appendix
\section{Formal residues and meromorphic germs}
\label{sec:formalresidue}

This appendix records the purely local residue calculus, which does not
require a contour or a standard-part domain. Let
\[
 \K((T))=\left\{\sum_{n\geq N}a_nT^n:N\in\Z,\ a_n\in\K\right\}.
\]
The negative part is finite. The positive part becomes a germ-analytic
function by \cref{thm:scale}; dividing by a power of $T$ gives an actual
meromorphic germ on a sufficiently small punctured neighborhood. Define
\[
 \Res_T\left(\sum_{n\geq N}a_nT^n\,dT\right)=a_{-1}.
\]

\begin{proposition}[Primitives and coefficient extraction]
A formal meromorphic differential $A(T)\,dT$ has a primitive in $\K((T))$
if and only if its residue is zero. Moreover, for $B(T)=\sum_{j\geq0}b_jT^j$,
\[
 \Res_T\frac{B(T)}{T^{n+1}}\,dT=b_n=\frac{B^{(n)}(0)}{n!}
 \qquad(n\geq0).
\]
\end{proposition}
\begin{proof}
A derivative has no $T^{-1}$ coefficient. Conversely, when $a_{-1}=0$,
\[
 \sum_{\substack{n\geq N\\n\neq-1}}\frac{a_n}{n+1}T^{n+1}
\]
is a well-defined Laurent series and differentiates to $A$. Coefficient
extraction in the second statement is immediate. The corresponding germ
identities follow from the local realization theorem.
\end{proof}

\begin{proposition}[Logarithmic residue]
For $0\neq f\in\K((T))$,
\[
 \Res_T\frac{f'(T)}{f(T)}\,dT=\ord_T f\in\Z.
\]
\end{proposition}
\begin{proof}
Write $f=T^m u$, where $m\in\Z$ and $u\in\K[[T]]$ is a unit. Then
$f'/f=m/T+u'/u$, and the second summand is a power series.
\end{proof}

\begin{theorem}[Change of coordinate and ramification]
Let $\psi(T)=T^m u(T)$, where $m\geq1$ and $u(0)\neq0$. For every
$A(Z)\in\K((Z))$,
\begin{equation}\label{eq:formal-change}
 \Res_T A(\psi(T))\psi'(T)\,dT
       =m\Res_Z A(Z)\,dZ.
\end{equation}
In particular, residues are invariant under every invertible local analytic
coordinate change.
\end{theorem}
\begin{proof}
For a Laurent monomial $A(Z)=Z^n$ with $n\neq-1$, the pullback is the
formal derivative of $\psi^{n+1}/(n+1)$ and has residue zero. For $n=-1$,
it is $\psi'/\psi$, whose residue is $m$ by the preceding proposition.
Only finitely many negative Laurent monomials occur in $A$, while a
nonnegative monomial cannot contribute a residue after substitution.
Thus coefficient extraction commutes with the substitution, and linearity
proves \eqref{eq:formal-change}. An invertible coordinate has $m=1$.
\end{proof}

\begin{theorem}[Rational residue theorem over $\K$]
For a rational differential $R(Z)\,dZ$ on the projective line over $\K$,
\[
 \sum_{a\in\K}\Res_a R(Z)\,dZ+\Res_\infty R(Z)\,dZ=0,
\]
where only finitely many terms are nonzero and
\[
 \Res_\infty R(Z)\,dZ
      =\Res_{T=0}\bigl(-T^{-2}R(T^{-1})\bigr)\,dT.
\]
\end{theorem}
\begin{proof}
Algebraic closedness gives a finite partial-fraction decomposition
\[
 R(Z)=Q(Z)+\sum_a\sum_{j=1}^{m_a}\frac{c_{a,j}}{(Z-a)^j}.
\]
The finite residue at $a$ is $c_{a,1}$. After substituting $Z=T^{-1}$,
the polynomial part has no $T^{-1}$ term. A summand with $j\geq2$ becomes
$-c_{a,j}T^{j-2}(1-aT)^{-j}$ and also has no $T^{-1}$ term. A summand
with $j=1$ has residue $-c_{a,1}$ at $T=0$. Summation proves the result.
\end{proof}

There is no inconsistency between this appendix and the global primitive
$L'(z)=1/z$ from \cref{thm:globallog}. That $L$ is locally analytic at each
\emph{nonzero} point; it is not a meromorphic Laurent germ at zero. A
primitive in the enormous category of locally analytic functions on a
punctured class-domain need not belong to $\K((T))$ at the puncture.

\section{Notation and conventions}

\renewcommand{\arraystretch}{1.16}
\begin{longtable}{@{}>{\raggedright\arraybackslash}p{0.27\textwidth}p{0.67\textwidth}@{}}
\toprule
Notation & Meaning\\
\midrule
\endfirsthead
\toprule
Notation & Meaning\\
\midrule
\endhead
$\No$, $\K=\No(i)$ & The surreal ordered field and its complexification.
Both are proper classes.\\
$\tmon^\gamma=\omega^{-\gamma}$ & Conway monomial notation, with the order
reversed so that positive exponents represent infinitesimal monomials.
It is not defined by $\exp(-\gamma\log\omega)$.\\
$v(z)$, $\lc(z)$ & Least normal-form exponent and its nonzero complex
coefficient; $v(0)=+\infty$.\\
$\OO$, $\mm$ & Finite valuation ring $\{v\geq0\}$ and infinitesimal
ideal $\{v>0\}$, respectively. In $\OO(U)$, the same conventional letter
instead denotes ordinary holomorphic functions on $U$.\\
$\st$, $U^\sharp$ & Standard-part map $\OO\to\C$ and the inverse image
$\st^{-1}(U)$.\\
$\GG_a$ & The algebra of Hahn-analytic germs at $a$, naturally isomorphic
to $\K[[T]]$ with $T=z-a$.\\
$\HH(U)$, $\vprof$ & Holomorphic profiles with one common well-ordered
set of exponents, and their least nonzero coefficient exponent.\\
$\hhat F$ & Canonical evaluation of a profile on $U^\sharp$ by complex
Taylor expansion in the infinitesimal displacement.\\
$\hint$, $\hoint$ & Coefficientwise Hahn-valued integration over an
ordinary complex contour. These symbols do not denote sharp limits of
Riemann sums.\\
$D_z$, $\partial$ & Variable differentiation and a strongly additive
derivation of the coefficient field, respectively.\\
$O(q^n)$ & In the formal worked examples, membership in the ideal
$q^n\C[[q]]$ (or its specified coefficient extension), not convergence
of a sequence in the full sharp topology.\\
\bottomrule
\end{longtable}
\renewcommand{\arraystretch}{1}

\begin{thebibliography}{99}
\addcontentsline{toc}{section}{References}
\markboth{References}{}

\bibitem{Alling}
N. L. Alling, \emph{Foundations of Analysis over Surreal Number Fields},
North-Holland Mathematics Studies \textbf{141}, North-Holland, 1987.
In particular, Chapter 7, ``Power Series: Formal and Hyper-Convergent.''

\bibitem{BMder}
A. Berarducci and V. Mantova, ``Surreal numbers, derivations and
transseries,'' \emph{Journal of the European Mathematical Society}
\textbf{20} (2018), 339--390.
\href{https://doi.org/10.4171/JEMS/769}{doi:10.4171/JEMS/769}.
Preprint: \url{https://arxiv.org/abs/1503.00315}.

\bibitem{BMgerms}
A. Berarducci and V. Mantova, ``Transseries as germs of surreal functions,''
\emph{Transactions of the American Mathematical Society}
\textbf{371} (2019), 3549--3592.
\href{https://doi.org/10.1090/tran/7428}{doi:10.1090/tran/7428}.
Preprint: \url{https://arxiv.org/abs/1703.01995}.

\bibitem{ConwayComplex}
J. B. Conway, \emph{Functions of One Complex Variable I}, second edition,
Graduate Texts in Mathematics \textbf{11}, Springer, 1978.
\href{https://doi.org/10.1007/978-1-4612-6313-5}{doi:10.1007/978-1-4612-6313-5}.

\bibitem{Gonshor}
H. Gonshor, \emph{An Introduction to the Theory of Surreal Numbers},
London Mathematical Society Lecture Note Series \textbf{110},
Cambridge University Press, 1986.

\bibitem{MS}
J. M\"uller and A. Strohmaier, ``The theory of Hahn-meromorphic functions,
a holomorphic Fredholm theorem, and its applications,''
\emph{Analysis \& PDE} \textbf{7} (2014), 745--770.
\href{https://doi.org/10.2140/apde.2014.7.745}{doi:10.2140/apde.2014.7.745}.
Preprint: \url{https://arxiv.org/abs/1205.0236}.

\bibitem{Neumann}
B. H. Neumann, ``On ordered division rings,''
\emph{Transactions of the American Mathematical Society}
\textbf{66} (1949), 202--252.

\bibitem{Orloff}
J. Orloff, \emph{18.04 Complex Variables with Applications},
MIT OpenCourseWare, Spring 2018, lecture notes, especially Topics 4--5
(Cauchy theory) and Topic 11 (argument principle).
\href{https://ocw.mit.edu/courses/18-04-complex-variables-with-applications-spring-2018/pages/lecture-notes/}{MIT OpenCourseWare lecture-note collection}.

\bibitem{RSS}
S. Rubinstein-Salzedo and A. Swaminathan, ``Analysis on surreal numbers,''
\emph{Journal of Logic and Analysis} \textbf{6}(5) (2014), 1--39.
Preprint and revised version: \url{https://arxiv.org/abs/1307.7392}.

\end{thebibliography}
\end{document}
